% concatenated from main.tex
\documentclass[11pt]{article}

\usepackage{amssymb}
\usepackage{amsthm}
\usepackage{amsmath}
\newtheorem{theorem}{Theorem}[section]
\newtheorem{remark}{Remark}[section]
\RequirePackage{algorithm}
\RequirePackage[noend]{algpseudocode}
\usepackage{xcolor}
\usepackage{silence}
\usepackage{hyperref}
\WarningFilter{latex}{Text page 7 contains only floats}

\newtheorem{axiom}{Axiom}
\newtheorem{claim}[axiom]{Claim}
\newtheorem{definition}{Definition}[section]
\newtheorem{prop}{Proposition}[section]
\newtheorem{lem}{Lemma}[section]


\def\R{\mathbb{R}}
\def\C{\mathbb{C}}
\def\Z{\mathbb{Z}}
\def\N{\mathbb{N}}
\def\Q{\mathbb{Q}}
\def\D{\mathbb{D}}
\def\Ex{\mathbb{E}}
\def\P{\mathbb{P}}
\def\T{\mathbb{T}}
\def\Sp{\mathbb{S}}
\def\hb{\hfil \break}
\def\ni{\noindent}
\def\x{{\xi}^{\ast}}
\def\i{\indent}
\def\a{\alpha}
\def\b{\beta}
\def\e{\epsilon}
\def\d{\delta}
\def\De{\Delta}
\def\g{\gamma}
\def\qq{\qquad}
\def\L{\Lambda}
\def\E{\cal E}
\def\G{\Gamma}
\def\F{\cal F}
\def\Hi{\cal H}
\def\K{\cal K}
\def\O{\cal O}
\def\A{\cal A}
\def\B{\cal B}
%\def\S{\cal S}
\def\L{\cal L}
\def\M{\cal M}
\def\N{\cal N}
%\def\P{\cal P}
\def\Om{\Omega}
\def\om{\omega}
\def\s{\sigma}
\def\th{\theta}
\def\Th{\Theta}
\def\z{\zeta}
\def\p{\phi}
\def\m{\mu}
\def\n{\nu}
\def\l{\lambda}
\def\Si{\Sigma}
\def\q{\quad}
\def\qq{\qquad}
\def\half{\frac{1}{2}}
\def\hb{\hfil \break}
\def\half{\frac{1}{2}}
\def\pa{\partial}
\def\ni{\noindent}
\def\i{\indent}
\def\half{{1 \over 2}}
\def\Om{\Omega}
\def\om{\omega}
%%% END DEFs


\begin{document}

% \begin{frontmatter}

% \title{Szeg\H{o} Theorem for Operator Orthogonal Polynomials}

% \author[inst2]{Nicholas H. Bingham}

% \cortext[cor1]{Corresponding author.}

% \author[inst1]{Badr Missaoui\corref{cor1}}

% \affiliation[inst2]{organization={Department of Mathematics, Imperial College London},%Department and Organization
%             addressline={e-mail: n.bingham@ic.ac.uk}, 
% }
% \affiliation[inst1]{organization={Moroccan Center For Game Theory, UM6P},%Department and Organization
%             e-mail={e-mail: badr.missaoui@um6p.ma}}

            	\renewcommand{\thefootnote}{\arabic{footnote}}
	%$^{,}$
	\begin{center}
		{\Large \textbf{ Szeg\H{o} Theorem for Operator Orthogonal Polynomials}} \\[0pt]
		~\\[0pt] Badr Missaoui \footnote[1]{Mohammed VI Polytechnic University, Morocco. E-mail: \texttt{badr.missaoui@um6p.com}}, Nicholas H. Bingham{\footnote[2]{Department of Mathematics, Imperial College London, UK.
				E-mail: \texttt{n.bingham@ic.ac.uk}}}
		\\[0pt]
		
		
		~\\[0pt]
	\end{center}



\begin{abstract}
We develop a Szeg\H{o} theory for operator-valued measures on the unit circle, extending the framework of matrix orthogonal polynomials to the infinite-dimensional setting. Our approach yields a new and direct proof of the operator Szeg\H{o} limit theorem based solely on orthogonal-polynomial theory.
\end{abstract}

{\bf Keyword}: Moment problem, orthogonal polynomials, Szeg\H{o} theory.
%\MSC[2020]{Primary 42C05; Secondary 42A10, 42A70}


\section{Introduction}
The classical Szeg\H{o} limit theorem, originally proven in \cite{Sze}, describes the asymptotic behavior of Toeplitz determinants and plays a fundamental role in the spectral theory of Toeplitz operators and orthogonal polynomials on the unit circle. Specifically, if $\mu$ is a probability measure on the unit circle such that $\log \mu' \in L^{1}$, then  
\begin{equation}
\lim_{n \rightarrow{ \infty}}\frac{\mathrm{det}(T_{n})}{\mathrm{det}(T_{n-1})}=\prod_0^{\infty} (I - {\alpha}_k {\alpha}_k^{\dag})
= \exp \int\ \log \mu'(\theta) d \theta /2 \pi ,
\end{equation}
where $T_n=[\mu_{i-j}]_{i,j=0}^n$ and $\{\mu_k\}_{k \in \mathbb{Z}}$ are the Fourier coefficients of $\mu$, $\alpha_k$ is the \textit{k-th} Verblunsky coefficient associated with $\mu$, and ${\alpha}_k^{\dag}$ is its Hermitian conjugate. This result, deeply connected to orthogonal polynomials, has been extensively generalized in various settings, including matrix-valued orthogonal polynomials \cite{DerHK,DelGK}.  

In this paper, we extend this framework to the setting of operator orthogonal polynomials, providing a Szeg\H{o}-type limit theorem in infinite-dimensional Hilbert spaces. The study of operator-valued polynomials is not merely of theoretical interest; it has significant implications in prediction theory, particularly in the analysis of continuous stochastic processes and large collections of time series \cite{Bos, DugBM}. The use of operator-valued functions and orthogonal polynomials in prediction theory dates back to the foundational works of Mandrekar and Salehi \cite{ManS} and Weron \cite{Wer}, who developed the theory of dilation and shift operators for nonstationary processes. A key concept in this area is the \textit{Kolmogorov decomposition}, which expresses a positive definite kernel as the inner product of two operator-valued sequences, enabling the analysis of nonstationary processes through dilation and shift operators.  \\

Operator orthogonal polynomials extend the concept of matrix-valued \textit{orthogonal} polynomials by considering polynomials whose coefficients are bounded operators on a Hilbert space rather than scalars or matrices. Formally, let $\mathcal{H}$ be a separable Hilbert space, and let $\mathcal{L}(\mathcal{H})$ denote the space of bounded linear operators on $\mathcal{H}$. Given a positive, self-adjoint operator-valued measure $\mu$, we consider a sequence of operator-valued polynomials $P_n(z)$, defined on the unit circle $\mathbb{T}$, with  
\[
P_n(z) = \sum_{k=0}^{n} a_k z^k, \quad a_k \in \mathcal{L}(\mathcal{H}).
\]  
These polynomials are said to be orthogonal with respect to $\mu$ if  
\begin{align*}  
\langle P_n(z), P_m(z) \rangle_{\mu} = 0, \quad \text{for } n \neq m,  
\end{align*}  
where the inner product is defined as  
\[
\langle f, g \rangle_{\mu} := \int_{\mathbb{T}} f(z)^{\dag} d\mu(z) g(z).
\]  
Here, $\mathbb{T}$ denotes the unit circle in the complex plane, and $\mu$ is a non-trivial operator-valued measure, i.e. one not supported on a finite set. The study of such polynomials is closely related to dilation theory, which seeks to extend a given operator to a larger space while preserving its essential properties. This theory, pioneered by Sz.-Nagy and Foia\c{s} \cite{SzNF} and further developed by Nikolskii \cite{Nik2}, has been applied to operator-valued functions by Gorniak and Weron \cite{GorW} and Weron \cite{Wer}.  \\

The generalization of the Szeg\H{o} limit theorem to operator settings have previously been established. Gohberg and Kaashoek \cite{GohK} provided an operator generalization using factorization arguments and Schur complement techniques, focusing on a class of positive block Toeplitz operators with Hilbert-Schmidt entries. Their results relied on the second regularized and Perelson determinants in place of the classical determinant. Subsequent research by B\"ottcher and Silbermann \cite{BotS} removed the self-adjointness requirement on the block Toeplitz matrices and relaxed the smoothness conditions imposed by Gohberg and Kaashoek. \\ 

The novelty of the present work lies instead in the method of proof and structural viewpoint. We provide a derivation the operator Szeg\H{o} limit theorem that relies exclusively on operator-valued orthogonal polynomials on the unit circle. Our method makes use of the moments of the operator-valued measure to represent operator orthogonal polynomials in terms of Schur complements (see, e.g., \cite{Zha}). The core of our method closely follows the matrix-valued results developed in Simon's monograph \cite[Section 2.13]{Sim2}, as well as in \cite{DamPS} and \cite{DerHK}, allowing a direct extension of the matrix theory to the operator setting. This leads naturally to operator analogues of several classical results, including:
\begin{itemize}
    \item Classical recursion relations for operator orthogonal polynomials,
    \item Operator versions of the kernel polynomials and Christoffel-Darboux formulas, and
    \item A new perspective on the Bernstein-Szeg\H{o} approximation in the operator setting.
    \item A new proof of the operator Szeg\H{o} limit theorem. Our decisive result (Theorems 7.1 and 7.2) provides an operator analogue of the classical Szeg\H{o} limit theorem.
\end{itemize}

The structure of this paper is as follows: Section 2 introduces the fundamental concepts and notation essential for our analysis of operator orthogonal polynomials. In Section 3, we define and establish the basic properties of operator-valued counterparts to matrix-valued orthogonal polynomials on the unit circle. Sections 4 through 6 develop the operator versions of the Verblunsky recurrence relations, the Christoffel-Darboux formulas, and the Bernstein-Szeg\H{o} approximation. Section 7 presents the proof of the operator Szeg\H{o} limit theorem within the framework of matrix orthogonal polynomials on the unit circle.

\section{Preliminaries} 
In this section, we study the fundamental properties of positive operator-valued measures, their spectral representations, and connection to stationary processes.\\

Let $\mathcal{B}(\mathbb{T})$ denote the class of all Borel measurable subsets of the unit circle $\mathbb{T}$. We define $\mathcal{L}(\mathcal{H},\mathcal{K})$ as the algebra of all bounded linear operators from $\mathcal{H}$ to $\mathcal{K}$ where $\mathcal{H}$ and $\mathcal{K}$ are separable spaces. When $\mathcal{H} = \mathcal{K}$, we abbreviate this to $\mathcal{L}(\mathcal{H})$. Recall that $A \in \mathcal{L}(\mathcal{H})$ is positive definite $(A>0)$ if $\langle A x, x \rangle >0$ for $0 \neq x \in \mathcal{H}$, and $A$ has a bounded inverse.

\subsection{Operator-Valued Measures}

A function $\mu: \mathcal{B}(\mathbb{T}) \to \mathcal{L}(\mathcal{H})$ is called a \textit{positive operator-valued measure} (or \textit{semi-spectral measure}) on $\mathbb{T}$ if, for each $ h \in \mathcal{H}$, the scalar measure
\[
\mu_h(\sigma) = \langle \mu(\sigma) h, h \rangle,
\]
is a positive Borel measure on $\mathbb{T}$.

It can be shown that for each $ j \in \mathbb{Z}$, there exists an operator $ \mu(j) \in \mathcal{L}(\mathcal{H}) $, known as the \textit{$j$th moment} of $\mu$, such that for each $ h \in \mathcal{H}$,
\[
\langle \mu(j) h, h \rangle = \int_{\mathbb{T}} e^{i j t} d\mu_h(t).
\]

We assume that $ \mu $ is \textit{absolutely continuous} with respect to the Lebesgue measure on $\mathbb{T}$. That is, there exists a weakly measurable function $ \mu' : \mathbb{T} \to \mathcal{L}(\mathcal{H}) $, such that $ \mu'(e^{it}) \geq 0 $ almost everywhere on $\mathbb{T}$, and for all $ h \in \mathcal{H} $,
\[
 d \langle \mu(t) h, h \rangle = \langle \mu'(e^{it}) h, h \rangle dt.
\]
The $ j $th moment of $ \mu $ then corresponds to the Fourier coefficient $ \widehat{\mu'}(j) $ of $ \mu' $, leading to the expansion
\[
 \mu'(e^{it}) = \sum_{j=-\infty}^{\infty} \mu(j) e^{ijt}.
\]

\subsection{Spectral Representation and Stationary Processes}

Let $ X = (X_t)_{t \in G} $ be a stationary stochastic process indexed by a locally compact abelian group $ G $, satisfying the covariance condition:
\[
 C(t, s) = \mathbb{E}(X_s^* X_t) = C'(t - s).
\]
By the \textit{classical spectral theorem} \cite{DunSa}, there exists a unique unitary representation $ U : G \to \mathcal{L}(\mathcal{H}) $ satisfying
\[
 X_t = U_t X_0.
\]
The operator $ U_t $ is called the \textit{shift operator}, which captures the evolution of the process over time. By \textit{Stone’s theorem} \cite{Sto}, the shift operator has a spectral representation:
\[
 U = \int_{\mathbb{T}} e^{i\lambda} E(d\lambda),
\]
where $ E $ is an \textit{operator-valued spectral measure} on $\mathbb{T}$. Consequently, the process $ X_t $ admits the integral representation
\[
 X_t = U^t X_0 = \int_{\mathbb{T}} e^{it\lambda} E(d\lambda) X_0 = \int_{\mathbb{T}} e^{it\lambda} \xi(d\lambda),
\]
where $ \xi(\Delta) = E(\Delta) X_0 $ is a measure-valued process.

If $ G = \mathbb{Z} $, the covariance function satisfies
\[
 C_n = \langle X_{m+n}, X_m \rangle = \int_{\mathbb{T}} e^{-in\lambda} \left\langle \xi(d\lambda), \xi(d\lambda) \right\rangle.
\]
By uniqueness of the Fourier transform, the spectral measure $ \mu $ satisfies:
\[
 \mu(\Delta) = X_0^* E(\Delta) X_0.
\]

\subsection{Kolmogorov Isomorphism and Prediction Theory}

A key result in stochastic prediction is the \textit{Kolmogorov Isomorphism Theorem}, which asserts the existence of a unitary equivalence between the process $ X_n $ and a weighted shift operator on $ L^2(\mu, \mathcal{H}) $. The correspondence
\[
 X_n \leftrightarrow e^{-in\cdot} I, \quad n\in \mathbb{Z},
\]
where $ I $ is the identity operator on $ \mathcal{H} $, reformulates the Kolmogorov-Wiener prediction problem as an approximation problem in $ L^2(\mu, \mathcal{H}) $.\\

From (\cite{ManS} Theorem 4.19), the space $ L^2(\mu, \mathcal{H}) $ consists of operator-valued functions satisfying
\[
 \text{tr}  \int_{\mathbb{T}} f^\dagger(x) \mu(dx) f(x)  < \infty,
\]
equipped with the inner products
\[
 \left\langle \left\langle f, g \right\rangle \right\rangle_R = \int_{\mathbb{T}} f^\dagger(x) \mu(dx) g(x), \quad \left\langle \left\langle f, g \right\rangle \right\rangle_L = \int_{\mathbb{T}} g(x) \mu(dx) f^\dagger(x).
\]

\noindent An operator-valued measure is called non-trivial if $$\text{tr} \left\langle \left\langle f, g \right\rangle \right\rangle_R >0,$$ for every non-zero polynomial $f$. Here, \text{tr} is the trace.

\noindent Moreover, 
$$\left\langle \left\langle f, g \right\rangle \right\rangle_\mathrm{L} = \left\langle \left\langle g, f \right\rangle \right\rangle_\mathrm{L}^{\dag}, ~~~
\left\langle \left\langle f, g \right\rangle \right\rangle_\mathrm{R} = \left\langle \left\langle g, f \right\rangle \right\rangle_\mathrm{R}^{\dag}$$
$$\left\langle \left\langle f^{*}, g^{*} \right\rangle \right\rangle_\mathrm{L} = \left\langle \left\langle g, f \right\rangle \right\rangle_\mathrm{R}^{\dag}, ~~~
\left\langle \left\langle f^{*}, g^{*} \right\rangle \right\rangle_\mathrm{R} = \left\langle \left\langle g, f \right\rangle \right\rangle_\mathrm{L}^{\dag}$$
$$\left\langle \left\langle f, \alpha g \right\rangle \right\rangle_\mathrm{L} = \alpha \left\langle \left\langle f, g \right\rangle \right\rangle_\mathrm{L}^{\dag}, ~~~
\left\langle \left\langle \alpha f, g \right\rangle \right\rangle_\mathrm{L} = \left\langle \left\langle f, g \right\rangle \right\rangle_\mathrm{L} \alpha ^{\dag}$$
$$\left\langle \left\langle f, g \alpha \right\rangle \right\rangle_\mathrm{R} = \left\langle \left\langle f, g \right\rangle \right\rangle_\mathrm{R}\alpha, ~~~
\left\langle \left\langle f, \alpha g \right\rangle \right\rangle_\mathrm{R} = \alpha^{\dag} \left\langle \left\langle f, g \right\rangle \right\rangle_\mathrm{R}.$$





Consider a family of Hilbert spaces \( \mathbf{H} = \{ \mathcal{H}_n \}_{n \in \mathbb{Z}} \), where each \( \mathcal{H}_n \) is a separable Hilbert space. A function \( C: \mathbb{Z} \times \mathbb{Z} \to \mathcal{L}(\mathcal{H}_j, \mathcal{H}_i) \) is called a \textit{positive definite kernel} if it satisfies the condition
\[
\sum_{i,j} \langle C_{i,j} h_j, h_i \rangle \geq 0,
\]
for all finitely supported sequences \( \{ h_n \}_{n \in \mathbb{Z}} \) in the Hilbert direct sum \( \bigoplus_{n \in \mathbb{Z}} \mathcal{H}_n \). \\

The following theorem establishes a fundamental link between shift operators and positive definite sequences of operators. For further background, see \cite{Con}.

\begin{theorem} \label{thm:kolmoDecomp}[Kolmogorov decomposition]
Let \( C \) be a positive definite kernel. Then there exist a Hilbert space \( \mathcal{K} \) and a collection of bounded operators \( V_n \in \mathcal{L}(\mathcal{H}_n, \mathcal{K}) \) such that:
\begin{itemize}
    \item \( C_{i,j} = V_i^* V_j \), for all \( i,j \in \mathbb{Z} \);
    \item \( \mathcal{K} = \bigvee_{n \in \mathbb{Z}} V_n \mathcal{H}_n \), where \( \bigvee \) denotes the closed linear span.
\end{itemize}
\end{theorem}

A particularly important case occurs when the family \( \mathbf{H} \) reduces to a single Hilbert space, i.e., when \( \mathcal{H}_n = \mathcal{H} \) for all \( n \in \mathbb{Z} \), and the positive definite kernel \( C \) satisfies the property:
\[
C_{i,j} = T_{j-i},
\]
for some function \( T: \mathbb{Z} \to \mathcal{L}(\mathcal{H}) \). In this case, \( C \) is referred to as a \textit{positive definite Toeplitz kernel}.

\begin{theorem} \label{thm:naimark}[Naimark dilation]
Let \( C \) be a positive definite Toeplitz kernel. Then there exist a Hilbert space \( \mathcal{K} \), a unitary operator \( S \in \mathcal{L}(\mathcal{K}) \), and a bounded operator \( Q \in \mathcal{L}(\mathcal{H}, \mathcal{K}) \) such that:
\begin{itemize}
    \item \( C_{i,j} = Q^* S^{j-i} Q \), for all \( i,j \in \mathbb{Z} \);
    \item \( \mathcal{K} = \bigvee_{n \in \mathbb{Z}} S^n Q \mathcal{H}\).
\end{itemize}
\end{theorem}

The shift operator, which plays a fundamental role in functional and stochastic analysis, is defined by the successive powers of the unitary dilation operator \( U \). It is important to note that Theorem (\ref{thm:naimark}) applies only to \textit{stationary} processes. For nonstationary processes, a more general result is available: the Naimark dilation extends to the Kolmogorov decomposition, as described in Theorem \ref{thm:kolmoDecomp}.

This Kolmogorov decomposition extends the Kolmogorov Isomorphism Theorem (KIT) \cite{Kol} to the setting of operator-valued non-stationary processes, providing a unique minimal dilation representation for the correlation function \( C \) of a stochastic process. This fundamental result, originally established by Mandrekar and Salehi \cite{ManS} (\S 6), builds on the pioneering work of Wiener and Masani \cite{WieM} and reinforces the perspective proposed by Masani \cite{Mas2}, which asserts a deep connection between dilation theory and shift operators.

Dilation theory offers a powerful framework for analyzing signals by shifting the focus from the signal itself to its associated dilation operator, much as Fourier analysis shifts the emphasis from the time-domain representation to the spectral domain. Unlike Fourier methods, which are generally limited to stationary processes, dilation theory does not impose stationarity constraints. 

In the case of stationary processes, dilation theory naturally leads to the Naimark dilation (Theorem \ref{thm:naimark}), providing a unitary dilation framework. However, for non-stationary processes, the appropriate generalization is given by the Kolmogorov decomposition (Theorem \ref{thm:kolmoDecomp}), which constructs a minimal isometric dilation that captures the underlying structure of the non-stationary process.



\section{Operator orthogonal polynomials}

In Szeg\H{o} theory, orthogonal polynomials play a prominent role. Here, we will need to work with operator-valued orthogonal polynomials with respect to operator-valued measures on the unit circle. \\

For an operator polynomial $P_{n}$ of degree $n$, we defined the \textit{reversed polynomial} $P_{n}^{*}$  of $P_{n}$ by $$
P^{*}_{n}(z)=z^{n}P_{n}(1/\bar{z})^{\dag}.
$$
\noindent We have $$(P^{*}_{n})^{*}= P_{n},$$ and for any $\alpha\in \mathcal{L}(\mathcal{H})$, $$(\alpha P_{n})^{*} = P_{n}^{*}\alpha^{\dag},~~(P_{n}\alpha )^{*} = \alpha^{\dag}P_{n}^{*}.$$

The following lemma will be a very useful characterization of positive definite 2 $\times$ 2 operator matrices.

\begin{lem}
The following are equivalent:
\begin{enumerate}
% [label=\roman*]
  \item The operator matrix 
  	   $
       \begin{pmatrix}
       A & B\\
       B^{*} & C
       \end{pmatrix}
       $
       is positive definite. \\
  \item $A>0$ and $C-B^{*}A^{-1}B>0$.
  \item $C>0$ and $A-B^{*}C^{-1}B>0$;
\end{enumerate}
here $C-B^{*}A^{-1}B>0$ is called the {\it Schur complement} of $A$.
\end{lem}

\begin{proof}
This follows immediately from the \it{Frobenius-Schur factorization} (\cite{Zha}), 
\begin{equation}\label{schurcomplement}
\begin{pmatrix}
A & B\\
B^{*} & C
\end{pmatrix}
=\begin{pmatrix}
I & 0\\
B^{*}A^{-1} & I
\end{pmatrix}
\begin{pmatrix}
A & 0\\
0 & C-B^{*}A^{-1}B
\end{pmatrix}
\begin{pmatrix}
I & A^{-1}B\\
0 & I
\end{pmatrix}.
\end{equation}
\end{proof}

Given an operator-valued positive measure on $\mathbb{T}$, we define its \textit{moments} for $j=-n, \cdots, n$ by $$\mu_{j}=\mu_{j}(d\mu)=\frac{1}{2\pi}\int_{-\pi}^{\pi} e^{-ij\theta} \mu(d\theta) ~~ \mathrm{and}~~ \mu_{-j}=\mu^{*}_{j}.$$
The right and left {\it Toeplitz operator matrices} $T^{R}_{n}$ and $T^{L}_{n}$ associated to $\mu$ are
\[ \label{toeplitz_TR}
T^{R}_{n}=
\begin{pmatrix}
\mu_{0} & \mu_{-1} & \cdots & \mu_{-n+1}\\
\mu_{1} & \mu_{0} & \cdots & \mu_{-n+2}\\
\vdots & \vdots & \ddots & \vdots\\
\mu_{n-1} & \mu_{n-2} & \cdots & \mu_{0}
\end{pmatrix},
%&
T^{L}_{n}=
\begin{pmatrix}
\mu_{0} & \mu_{1} & \cdots & \mu_{n-1}\\
\mu_{-1} & \mu_{0} & \cdots & \mu_{n-2}\\
\vdots & \vdots & \ddots & \vdots\\
\mu_{-n+1} & \mu_{-n+2} & \cdots & \mu_{0}
\end{pmatrix}.
\]

\noindent We also have 
\begin{align*}
T^{R}_{n+1}&=
\begin{pmatrix}
T^{R}_{n} & \nu_{n}\\
\nu^{*}_{n} & \mu_{0}
\end{pmatrix},
&
T^{L}_{n+1}&=
\begin{pmatrix}
T^{L}_{n} & \xi_{n}\\
\xi^{*}_{n} & \mu_{0}
\end{pmatrix},
\end{align*}

\noindent where  
\begin{align*}
\nu_{n}&=
\begin{pmatrix}
\mu_{-n}\\
\mu_{-n+1}\\
\vdots \\
\mu_{-1}
\end{pmatrix},
&
\xi_{n}&=
\begin{pmatrix}
\mu_{n}\\
\mu_{n-1}\\
\vdots \\
\mu_{1}
\end{pmatrix},
\end{align*}

\noindent and we define the {\it Schur complements} $\kappa_{n}^{R,L}$ of $\mu_{0}$ {\it in} $T^{R,L}_{n+1}$ as
\begin{eqnarray*}
\kappa_{n}^{R} &=& \mathrm{SC}(T^{R}_{n+1}) = \mu_{0} - \nu^{*}_{n} T^{-R}_{n}\nu_{n}, \\
\kappa_{n}^{L} &=& \mathrm{SC}(T^{L}_{n+1}) = \mu_{0} - \xi^{*}_{n} T^{-L}_{n}\xi_{n},
\end{eqnarray*}
\noindent where  $T^{-R,-L}_{n}= (T^{R,L}_{n})^{-1}$.

\begin{remark} 
If $P(z)=\sum_{k=0}^{n}p_k z^k$ and $Q(z)=\sum_{k=0}^{n}q_k z^k$, then $$\left\langle \left\langle P, Q \right\rangle \right\rangle_\mathrm{R}= \sum_{k,j=0}^{n}p_k^{\dag} T_{k-j}^{R}q_j,\quad \left\langle \left\langle P, Q \right\rangle \right\rangle_\mathrm{L}= \sum_{k,j=0}^{n}q_j T_{k-j}^{L}p_j^{\dag}.$$
This shows that the positivity of $\mu$ implies the positivity of $T^{R,L}_{n}$ and $\kappa_{n}^{R,L}$, hence they have bounded inverses. Indeed, we have 
\begin{eqnarray*}
0 &<& \int_{\mathbb{T}}\left\langle \left\langle d\mu(\sigma) \sum_{j=0}^{n}p_j z^j, \sum_{k=0}^{n}p_k z^k \right\rangle \right\rangle_\mathrm{R,L} \\
  &<& \sum_{j,k=0}^{n}\int_{\mathbb{T}}z^{j-k}\left\langle \left\langle d\mu(\sigma) p_j, p_k \right\rangle \right\rangle_\mathrm{R,L}=\sum_{j,k=0}^{n}\left\langle \left\langle T_{k-j} p_j, p_k \right\rangle \right\rangle_\mathrm{R,L},
\end{eqnarray*}
so $T^{R,L}_{n}$ is positive definite, and has a bounded inverse $T^{-R,-L}_{n}$.  
\end{remark} 

\section{Verblunsky recursion}

In this section, we show that the right and left orthogonal polynomials follow the classical Verblunsky recurrence relations.\\
\noindent In our approach, we will define the monic operator-valued polynomials in terms of the Schur complements. The notion of Schur complements (SC) is relatively simple but suprisingly strong.

\begin{prop}
Monic polynomials such that 

\begin{eqnarray*}
\Phi_{n}^{R}(z)&=&\mathrm{SC} 
\begin{pmatrix}
\mu_{0} & \mu_{-1} & \cdots & \mu_{-n+1} & \mu_{-n}\\
\mu_{1} & \mu_{0} & \cdots & \mu_{-n+2} & \mu_{-n+1}\\
\vdots & \vdots & \ddots & \vdots & \vdots\\
\mu_{n-1} & \mu_{n-2} & \cdots & \mu_{0}  & \mu_{-1}\\
I & zI & \cdots & z^{n-1}I & z^{n}I
\end{pmatrix} \\
&=& z^{n}I-[I~~zI~~ \cdots ~~ z^{n-1}I]~T^{-R}_{n}
\begin{pmatrix}
\mu_{-n}\\
\mu_{-n+1}\\
\vdots \\
\mu_{-1}
\end{pmatrix},
\end{eqnarray*}

\begin{eqnarray*}
\Phi_{n}^{L}(z)&=&\mathrm{SC} 
\begin{pmatrix}
\mu_{0} & \mu_{-1} & \cdots & \mu_{n-1} & I\\
\mu_{1} & \mu_{0} & \cdots & \mu_{n-2} & zI\\
\vdots & \vdots & \ddots & \vdots & \vdots\\
\mu_{n-1} & \mu_{n-2} & \cdots & \mu_{0}  & z^{n-1}I\\
\mu_{-n} & \mu_{n-1} & \cdots & \mu_{-1} & z^{n}I
\end{pmatrix} \\
&=& z^{n}I-[\mu_{-n} ~~ \mu_{-n+1} ~~ \cdots ~~ \mu_{-1}]~T^{-L}_{n}
\begin{pmatrix}
I\\
zI\\
\vdots \\
z^{n-1}I
\end{pmatrix}
\end{eqnarray*}
are orthogonal, i.e., for any $k,j \geq 0$,
$$
\left\langle \left\langle \Phi_{k}^{R},\Phi_{j}^{R}\right\rangle\right\rangle_{\mathrm{R}} = \delta_{kj}\kappa_{k}^{R},  \qquad
\left\langle \left\langle \Phi_{k}^{L},\Phi_{j}^{L}\right\rangle\right\rangle_{\mathrm{L}} =\delta_{kj}\kappa_{k}^{L}.
$$

\end{prop}
\begin{proof}
We prove just the first relation; the second is similar.\\
\noindent For any $0 \leq m \leq n-1$ and $z=e^{i\theta}$
\begin{eqnarray*}
\left\langle \left\langle z^{m}I,\Phi_{n}^{R}\right\rangle\right\rangle_{\mathrm{R}} &=& 
\int_{-\pi}^{\pi}e^{-im\theta} \mu(\theta) 
(e^{i n \theta} - 
[I~~e^{i\theta}~~\cdots~~e^{i(n-1)\theta}]T^{-R}_{n}\nu_{n})d\theta\\
&=& \mu_{m-n} - [\mu_{m}~~\cdots~~\mu_{m-n+1}]T^{-R}_{n}\nu_{n}\\
&=& \mu_{m-n} - \mu_{m-n} = 0.
\end{eqnarray*}
If $m=n$, then 
\begin{eqnarray*}
\left\langle \left\langle \Phi_{n}^{R},\Phi_{n}^{R}\right\rangle\right\rangle_{\mathrm{R}} &=& 
\mu_{0} - [\mu_{n}~~\cdots~~\mu_{1}]T^{-R}_{n}\nu_{n} = \kappa_{n}^{R}.
\end{eqnarray*}
\end{proof}

\begin{theorem}
The monic orthogonal polynomials $\Phi^{R}_{n}$ and $\Phi^{L}_{n}$ obey the following recursion relations:
\begin{eqnarray} \label{monic_recursion}
\Phi_{n+1}^{R} &=& z \Phi_{n}^{R} + \Phi_{n}^{L,*} \Phi_{n+1}^{R}(0), \\
\Phi_{n+1}^{L} &=& z \Phi_{n}^{L} + \Phi_{n+1}^{L}(0) \Phi_{n}^{R,*}.
\end{eqnarray} 
\end{theorem}     
\begin{proof}
For the first recursion, observe that for any $1 \leq i \leq n$, we have 
$$\left\langle \left\langle \Phi_{n+1}^{R} - z \Phi_{n}^{R}, z^{i}I\right\rangle\right\rangle_{\mathrm{R}}= \left\langle \left\langle \Phi_{n}^{L,*}, z^{i}I\right\rangle\right\rangle_{\mathrm{R}} = 0,$$
and so $\Phi_{n+1}^{R} - z \Phi_{n}^{R}$ and $\Phi_{n}^{L,*}$ are proportional. Setting $z=0$ gives the constant of proportionality as $\Phi_{n+1}^{R}(0)$. 
 Similarly for the other claim.
 \end{proof} 
 
\noindent Based on Miamee and Salehi \cite{MiaS1}, the operator $\kappa^{R,L}_{n}$  has a well-defined square root that is positive definite. Therefore, we define the normalized orthogonal polynomials as follows:
$$
\varphi_{n}^{R} =\Phi_{n}^{R} \kappa^{-R/2}_{n}~~~~\mathrm{and}~~~~\varphi_{n}^{L} =\kappa^{-L/2}_{n} \Phi_{n}^{L}.
$$
\noindent One easily verifies 
$$
\left\langle \left\langle \varphi_{n}^{R,L},\varphi_{n}^{R,L} \right\rangle\right\rangle_{\mathrm{R,L}}=\kappa^{-R,L/2}_{n}\left\langle \left\langle \Phi_{n}^{R,L},\Phi_{n}^{R,L} \right\rangle\right\rangle_{\mathrm{R,L}}\kappa^{-R,L/2}_{n}=I.
$$

\noindent Defining 
$$
\rho_{n}^{R}=\kappa_{n+1}^{R/2}\kappa_{n}^{-R/2} ~~\mathrm{and}~~ \rho_{n}^{L}=\kappa_{n}^{-L/2}\kappa_{n+1}^{L/2},
$$
\noindent one can easily show 
\begin{eqnarray}\label{szego-recurrence}
z \varphi_{n}^{R} - \varphi_{n+1}^{R} \rho^{R}_{n}&=&  \varphi_{n}^{L,*}(\alpha_{n}^{R})^{\dag},  \label{szego-recurrence_right} \\
z \varphi_{n}^{L} - \rho^{L}_{n} \varphi_{n+1}^{L}&=&  (\alpha_{n}^{L})^{\dag} \varphi_{n}^{R,*}, \label{szego-recurrence_left}
\end{eqnarray} 
\noindent where 
\begin{eqnarray}
\alpha_{n}^{R} &=& -(\kappa^{-R/2}_{n})^{\dag} \Phi_{n+1}^{R}(0)^{\dag} \kappa^{L/2}_{n},  \label{eqn:alph_R}\\
\alpha_{n}^{L} &=& -\kappa^{R/2}_{n} \Phi_{n+1}^{L}(0)^{\dag} (\kappa^{-L/2}_{n})^{\dag}.
\end{eqnarray}

\noindent One also has 
$$
\kappa^{-L}_{n}=(\rho_{n-1}^{L} \dots \rho_{0}^{L})^{2}~~\mathrm{and}~~\kappa^{-R}_{n}=(\rho_{0}^{R} \dots \rho_{n-1}^{R})^{2}.
$$

\begin{prop}\label{prop4.2}
\hfill
\begin{enumerate}
    \item The operators $\alpha_{n}^{R}$ and $\alpha_{n}^{L}$ are equal, denoted $\alpha_n$.
    \item $\rho^{L}_{n} = (I - \alpha_{n}^{\dag}\alpha_{n})^{1/2}$ and $\rho^{R}_{n} = (I - \alpha_{n}\alpha_{n}^{\dag})^{1/2}$.
\end{enumerate}
\end{prop}

\begin{proof}
\noindent 
1) Observing that $\left\langle \left\langle \varphi_{n}^{R},\varphi_{n}^{L,*}\right\rangle\right\rangle_{\mathrm{R}} = \left\langle \left\langle \varphi_{n}^{R,*},\varphi_{n}^{L} \right\rangle\right\rangle_{\mathrm{L}} $, and using equations (\ref{szego-recurrence}), we can easily prove that $\alpha_{n}^{R}=\alpha_{n}^{L}$.

2) From (\ref{szego-recurrence_left}), we have 

\begin{eqnarray*}
||z\varphi_{n}^{L}||_{\mathrm{L}}- (\rho^{L}_{n})^{2} &=& \alpha_{n}^{\dag} \left\langle \left\langle \varphi_{n}^{R,*},\varphi_{n}^{R,*} \right\rangle\right\rangle_{\mathrm{L}} \alpha_{n}\\
&=& \alpha_{n}^{\dag}\alpha_{n}.
\end{eqnarray*}
Since $||z\varphi_{n}^{L}||_{\mathrm{L}} =1$, we have  $(\rho^{L}_{n})^{2} = 1 - \alpha_{n}^{\dag}\alpha_{n}.$\\
For $\rho^{R}_{n}$, the proof is similar.
\end{proof}

\noindent Using Proposition \ref{prop4.2}, we write,  
\begin{equation} \label{eq:recurrence_matrix}
      \begin{pmatrix}
        \varphi_{n}^{L}(z) \\
        \varphi_{n}^{R,*}(z) 
        \end{pmatrix}= A^{L}(\alpha_{n},z)\begin{pmatrix}
        \varphi_{n}^{L}(z) \\
        \varphi_{n}^{R,*}(z) 
      \end{pmatrix}
\end{equation}

where
\begin{equation} \label{eq:recurrence_inv_matrix}
A^{L}(\alpha,z) = \begin{pmatrix}
z(\rho^{L})^{-1} & -(\rho^{L})^{-1} \alpha^{\dag}\\
-z(\rho^{R})^{-1} \alpha & (\rho^{R})^{-1}
\end{pmatrix}.
\end{equation}


\section{Christoffel-Darboux Formulae}

This section introduces operator-valued right and left kernel polynomials and uses the recurrence formulae to derive the Christoffel-Darboux formulae. 

\begin{prop}\label{prop_kernelCD}
Define the right and left kernel polynomials of degree $n$ as
$$K_{n}^{R}(w,z) = \sum_{k=0}^{n}\varphi_{k}^{R}(z)\varphi_{k}^{R}(w)^{\dag} ,\quad K_{n}^{L}(w,z) = \sum_{k=0}^{n}\varphi_{k}^{L}(z)^{\dag}\varphi_{k}^{L}(w),$$
Then, we have
\begin{enumerate}
%[label=(\roman*)]
\item
$
K_{n}^{R}(w,z)=[I ~ zI ~\cdots ~ z^{n}I]~T_{n+1}^{-R}~\begin{bmatrix}
I \\
w^{-1}I \\
\vdots \\
w^{-n}I
\end{bmatrix}.
$
\item 
$
K_{n}^{L}(w,z)=[I ~ zI ~\cdots ~ z^{-n}I]~T_{n+1}^{-L}~\begin{bmatrix}
I \\
w^{1}I \\
\vdots \\
w^{n}I
\end{bmatrix}.
$
\item \begin{align} \label{cd_right}
K_{n}^{L}(w,z) &=&\frac{\varphi_{n}^{R,*}(w)^{\dag}\varphi_{n}^{R,*}(z) - \bar{w}z\varphi_{n}^{L}(w)^{\dag}\varphi_{n}^{L}(z)}{1-\bar{w}z}, \\ 
\label{cd_left}
K_{n}^{R}(w,z) &=&\frac{\varphi_{n}^{L,*}(w)\varphi_{n}^{L,*}(z)^{\dag} - \bar{w}z\varphi_{n}^{R}(w)\varphi_{n}^{R}(z)^{\dag}}{1-\bar{w}z}. 
\end{align}
\end{enumerate}
\end{prop}

\begin{proof}
We write $$
K_{n}^{R}(w,z)=[Z ~ z^{n}I]~T_{n+1}^{-R}~\begin{bmatrix}
W^{*}\\
w^{-n}I
\end{bmatrix},
$$ with $Z=[I ~ zI ~\cdots ~ z^{n-1}I]$ and $W=[I ~ wI ~\cdots ~ w^{n-1}I]$.\\
Writing  
$$
T^{R}_{n+1}=
\begin{pmatrix}
T^{R}_{n} & \nu_{n}\\
\nu^{*}_{n} & \mu_{0}
\end{pmatrix}
$$ 
and using (\ref{schurcomplement}), we have 
$$
T^{-R}_{n+1}=
\begin{pmatrix}
A & \gamma\\
\gamma^{*} & \alpha
\end{pmatrix},
$$
where 
$$
A = T^{-R}_{n} + T^{-R}_{n}\nu_{n}\kappa_{n}^{-R}\nu_{n}^{*}T^{-R}_{n},~~~ \gamma = -T^{-R}_{n}\nu_{n}\kappa_{n}^{-R},~~~ \alpha = \kappa_{n}^{-R}.
$$
Using the equality 
$$
[I ~ zI ~\cdots ~ z^{n-1}I]T^{-R}_{n}\nu_{n}\kappa_{n}^{-R/2}
= z^{n}\kappa_{n}^{-R/2} - \varphi_{n}^{R}(z),
$$
we rewrite $K_{n}^{R}(w,z)$ as follows:
\begin{align*}
K_{n}^{R}(w,z)
&= ZAW^{*} + z^{n}\gamma^{*}X^{*}+Z\gamma w^{-n} + z^{n}\alpha x^{-n}\\
&= Z \Bigl(T_{n}^{-R} \;+\; T_{n}^{-R}\,\nu_{n}\,\kappa_{n}^{-R}\,\nu_{n}^{*}\,T_{n}^{-R}\Bigr) W^{*}
   \;-\; z^{n}\,\kappa_{n}^{-R}\,\nu_{n}^{*}\,T_{n}^{-R}\\
   \;& \quad -\; Z\,T_{n}^{-R}\,\nu_{n}\,\kappa_{n}^{-R}\,w^{-n}   
   \;+\; z^{n}\,\kappa_{n}^{-R} w^{-n}\, \\[6pt]
&= Z\,T_{n}^{-R}\,W^{*}
   \;+\;\bigl(z^{n}\,\kappa_{n}^{-R/2}\;-\;\varphi_{n}^{R}(z))\,
    \bigl(w^{-n}\,\kappa_{n}^{-R/2}\;-\;\varphi_{n}^R(w)^{\dag}) \\
&\quad -\;w^{-n}\,\bigl(z^{n}\,\kappa_{n}^{-R/2}\;-\;\varphi_{n}^{R}(z)\bigr)\,\kappa_{n}^{-R/2}
   \;-\;z^{n}\,\kappa_{n}^{-R/2}\,\bigl(w^{-n}\,\kappa_{n}^{-R/2}\;-\;\varphi_{n}^{R}(w)^{\dag}) \\
&\quad +\;z^{n}\,w^{-n}\,\kappa_{n}^{-R} \\[6pt]
&= K_{n-1}^{R}(w,z) + \varphi_{n}^{R}(z) \varphi_{n}^{R}(w)^{\dag}.
\end{align*}
Summing, the telescoping sum gives 1. as required.\\

\noindent For 3.
$$
F_{n}^{L}(z)=\begin{pmatrix}
\varphi_{n}^{L}(z) \\
\varphi_{n}^{R,*}(z) 
\end{pmatrix}, ~~ 
J=\begin{pmatrix}
I & 0 \\
0 & -I
\end{pmatrix}, ~~
\tilde{J}=\begin{pmatrix}
\bar{w}zI & 0 \\
0 & -I
\end{pmatrix}.
$$
Then 
$$
F_{n+1}^{L}(z)= A^{L}(\alpha_{n},z)F_{n}^{L}(z)
$$ 
and 
$$
A^{L}(\alpha_{n},w)^{\dag} J A^{L}(\alpha_{n},z)=\begin{pmatrix}
\bar{w}zI & 0 \\
0 & -I
\end{pmatrix} = \tilde{J}.
$$ 
Thus 
$$
F_{n+1}^{L}(w)^{\dag} J F_{n+1}^{L}(z) = F_{n}^{L}(w)^{\dag} A^{L}(\alpha_{n},w)^{\dag} J A^{L}(\alpha_{n},z) F_{n}^{L}(z) = F_{n}^{L}(w)^{\dag} \tilde{J} F_{n}^{L}(z)
$$ 
and hence
$$
\varphi_{n+1}^{L}(w)^{\dag}\varphi_{n+1}^{L}(z) - \varphi_{n+1}^{R,*}(w)^{\dag}\varphi_{n+1}^{R,*}(z)=\bar{w}z\varphi_{n}^{L}(w)^{\dag}\varphi_{n}^{L}(z) - \varphi_{n}^{R,*}(w)^{\dag}\varphi_{n}^{R,*}(z),
$$ 
which gives the second part of (a). Now, if we let 
$$
Q_{n}^{L}(z,w)=\varphi_{n+1}^{R,*}(w)^{\dag}\varphi_{n+1}^{R,*}(z) - \varphi_{n+1}^{L}(z)^{\dag}\varphi_{n+1}^{L}(z),
$$ 
we see 
\begin{eqnarray*}
Q_{n+1}^{L}(z,w) - Q_{n}^{L}(z,w) &=& \varphi_{n}^{R,*}(w)^{\dag}\varphi_{n}^{R,*}(z) - \bar{w}z\varphi_{n}^{L}(w)^{\dag}\varphi_{n}^{L}(z) \\
& &- \varphi_{n}^{R,*}(w)^{\dag}\varphi_{n}^{R,*}(z) + \varphi_{n}^{L}(w)^{\dag}\varphi_{n}^{L}(z)\\
&=&(1-\bar{w}z)\varphi_{n}^{L}(w)^{\dag}\varphi_{n}^{L}(z).
\end{eqnarray*}
Summing over $n$ the proof for (\ref{cd_right}) follows since $Q_{-1}^{L}(z,w)=0$.\\
The proof for (\ref{cd_left}) is similar.
\end{proof}

\section{Bernstein-Szeg\H{o} Approximation}
In this section, we will reconstruct the operator-valued measure from the operator orthogonal polynomials via the Bernstein–Szeg\H{o} approximation.\\

\noindent Given a nontrivial operator-valued measure $d\mu$, with Verblunsky coefficients $\{\alpha_{n}\}_{n=0}^{\infty}$, we will define measures $d\mu^{(n)}$ with
\[
\alpha_{j}(d\mu^{(n)})= \left\{ 
	\begin{array}{ll}
        \alpha_{j},& \mbox{ $j\leq n$}  \\
        \mathbf{0},  & \mbox{ $j\geq n+1$}\\
    \end{array}
    \right.
\]
\noindent then $d\mu^{(n)} \rightarrow d\mu$ weakly (as in Simon \cite{Sim2} in the scalar case). As \cite{Sim2}, the proof below of the Szeg\H{o} theorem hinges on this.\\

\noindent As a preliminary, we need the following theorem:

\begin{theorem} \label{bernstein_zeros} We have
\begin{enumerate}
% [label=(\roman*)]
\item For $z \in \mathbb{T}$, all of $\varphi_{n}^{R,*}(z)$, $\varphi_{n}^{L,*}(z)$, $\varphi_{n}^{R}(z)$, $\varphi_{n}^{L}(z)$ are invertible.
\item For any $z \in \mathbb{T}$, $\varphi_{n}^{R}(z)\varphi_{n}^{R}(z)^{\dag} = \varphi_{n}^{L}(z)^{\dag}\varphi_{n}^{L}(z).$
\end{enumerate}

\end{theorem}
\begin{proof}
For 1., assume that $0< |w| \leq 1$. If there exists $c \neq 0$ such that $\varphi_{n}^{R}(w)^{\dag}c=0$, we also have that $\varphi_{n}^{L,*}(w)^{\dag}c=0$.\\
\noindent Then, propositions \ref{prop_kernelCD} lead to
$$
[I ~ zI ~\cdots ~ z^{n}I]~T_{n+1}^{-R}~\begin{bmatrix}
c \\
w^{-1}c \\
\vdots \\
w^{-n}c
\end{bmatrix}=0,~~ \mathrm{for}~~z \in \mathbb{C},
$$
which contradicts the invertibility of $T_{n+1}^{-R}$, unless $c=0$. This proves 1. for $\varphi_{n}^{L,*}(z)$ and $\varphi_{n}^{R}(z)$. If $z = e^{i\theta}$, $\varphi_{n}^{R,*}(z)= e^{-in\theta} \varphi_{n}^{R}(z)^{\dag}$ and $\varphi_{n}^{L}(z)= e^{in\theta} \varphi_{n}^{L,*}(z)^{\dag}$ are also invertible. \\
\indent For 2.,  put $z=w$ in equation (\ref{cd_right}).  One has $\varphi_{n}^{R,*}(z)^{\dag}\varphi_{n}^{R,*}(z) = \varphi_{n}^{L}(z)^{\dag}\varphi_{n}^{L}(z)$, which can be rewritten as 
$$
\varphi_{n}^{R}(z)\varphi_{n}^{R}(z)^{\dag} = \varphi_{n}^{L}(z)^{\dag}\varphi_{n}^{L}(z),
$$
as required.
% \indent For 2., Let $\varphi_{n}^{R}(z_{0}) = 0$ and define $p_{n-1}$ such that $(z-z_{0})p_{n-1}= \varphi_{n}^{R}$. The polynomial $p_{n-1}$ is of degree $n-1$, which implies $\left\langle \left\langle p_{n-1},\varphi_{n}^{R} \right\rangle\right\rangle_{\mathrm{R}} = 0$. So 
% $$
% \left\langle \left\langle zp_{n-1}, zp_{n-1}\right\rangle\right\rangle_{\mathrm{R}}
% =  \left\langle \left\langle p_{n-1}, p_{n-1}\right\rangle\right\rangle_{\mathrm{R}}
% = |z_{0}|^{2}\left\langle \left\langle p_{n-1}, p_{n-1}\right\rangle\right\rangle_{\mathrm{R}}+ 
% \left\langle \left\langle \varphi_{n}^{R}, \varphi_{n}^{R}\right\rangle\right\rangle_{\mathrm{R}},$$ or $$(1-|z_{0}|^{2})\left\langle \left\langle p_{n-1}, p_{n-1}\right\rangle\right\rangle_{\mathrm{R}} = \left\langle \left\langle \varphi_{n}^{R}, \varphi_{n}^{R}\right\rangle\right\rangle_{\mathrm{R}},
% $$ 
% from which we conclude $|z_0| < 1$, that is, the zeros of $\varphi_{n}$ lie in $\mathbb{D}$. Since $\varphi_{n}^{R,*}(z_{0})=0$ if and only if $\varphi_{n}^{R}(1/\bar{z}_{0})=0$, the zeros of $\varphi_{n}^{R,*}$ lie in $\mathbb{C}\backslash \mathbb{\bar{D}}$.
\end{proof}

As in the scalar case, we can use the invertibility in Theorem (\ref{bernstein_zeros}) to write 
\begin{equation} \label{eq:BS_measure}
d\mu^{(n)}(\theta) = [\varphi_{n}^{R}(e^{i\theta})\varphi_{n}^{R}(e^{i\theta})^{\dag}]^{-1}\frac{d\theta}{2\pi} = \mu'^{(n)}(\theta)\frac{d\theta}{2\pi}.
\end{equation} 
Also, directly from the definition of the right orthogonal polynomials, we have the right Bernstein-Szeg\H{o} approximation to $\mu$:
\[
d\mu^{(n)}(\theta) = [\varphi_{n}^{R,*}(e^{i\theta})^{\dag}\varphi_{n}^{R,*}(e^{i\theta})]^{-1}\frac{d\theta}{2\pi};
\]
\noindent the corresponding left Berstein-Szeg\H{o} approximation of ${\mu}_n     $ to $\mu$ is 
\[
d\mu^{(n)}(\theta) = [\varphi_{n}^{L,*}(e^{i\theta})\varphi_{n}^{L,*}(e^{i\theta})^{\dag}]^{-1}\frac{d\theta}{2\pi}.
\]

\begin{theorem}
The operator-valued measure $d\mu^{(n)}$ is normalized and its right operator orthogonal polynomials for $j=0, \cdots , n$ are $\{\varphi_{j}^{R}\}_{j=0}^{n}$, and for $j>n$, 
\begin{equation}\label{eq1_bernstein}
    \varphi_{j}^{R}(z; d\mu^{(n)})=z^{j-n}\varphi_{n}^{R}(z; d\mu).
\end{equation}

\noindent The Verblunsky coefficients for $d\mu^{(n)}$ are 
\begin{eqnarray}\label{eq:BD_approx_bis}
\alpha_{j}(d\mu^{(n)})= \left\{ 
	\begin{array}{ll}
        \alpha_{j}(d\mu),& \mbox{ $j\leq n$} , \\
        \mathbf{0},  & \mbox{ $j\geq n+1$}.\\
    \end{array}
    \right.    
\end{eqnarray} 
\end{theorem}
\begin{proof}
Let $\langle\langle \cdot, \cdot\rangle \rangle_{R}$ be the inner product associated with $\mu^{(n)}$. By direct calculation, 
\begin{equation}
\langle\langle \varphi_{n}^{R}, \varphi_{n}^{R} \rangle \rangle_{R} = \mathbf{1},
\end{equation}
and for $j = 0,1, \cdots, n-1$,
\begin{equation}
\langle\langle \varphi_{n}^{R}, \varphi_{j}^{R} \rangle \rangle_{R} = 0.
\end{equation}
So the family $\{\varphi_{n}^{R}\}_{j=0}^{n}$ is an orthonormal basis with respect to the inner product $\langle\langle \cdot, \cdot\rangle \rangle_{R}$. Also,
\begin{eqnarray*}
    \langle\langle e^{ij\theta}, \varphi_{n}^{R} \rangle \rangle_{R} &=&\frac{1}{2\pi}\int_{0}^{2\pi}e^{ij\theta}(\varphi_{n}^{R}(e^{i\theta})^{\dag})^{-1}d\theta\\
    &=& \frac{1}{2\pi}\oint e^{i(n-j-1)\theta}(\varphi_{n}^{R, *}(e^{i\theta}))^{-1}d\theta = 0,
\end{eqnarray*}
\noindent since $n-k-1 \geq 0$ and $\varphi_{n}^{R, *}(e^{i\theta})^{-1}$ is analytic in a neighborhood of $\bar{\mathbb{D}}$ by Theorem (\ref{bernstein_zeros}). This proves $\varphi_{n}^{R}$ is an OPUC for $d\mu_{n}$ and (\ref{eq1_bernstein}) holds. \\
By (\ref{eqn:alph_R}), if $\varphi^{R}_{k+1}(0)=0$, then $\alpha_k=0$, then by (\ref{eq1_bernstein}) $\varphi_{n+j}^{R}(0; d\mu^{(n)})=0$, which implies (\ref{eq:BD_approx_bis}).
\end{proof}
\noindent Next, measures with the properties above are essentially unique:
\begin{theorem} \label{geronimus_thm}
    Let $d\mu$ and $d\nu$ be two nontrivial operator-valued measures on $\mathbb{T}$ such that for $n$, 
    \begin{equation}\label{eq:Phi_N}
           \varphi_{n}^{R,L}(z;d\mu) = \varphi_{n}^{R,L}(z;d\nu).
    \end{equation} 
Then 
    \begin{eqnarray} \label{eq:Geronimus_op}
    	\begin{array}{ll}
            (i)~~\varphi^{R,L}_{j}(z;d\mu) = \varphi_{j}^{R,L}(z;d\nu)& \mbox{   $j=0,1,\cdots,n-1$} , \\
            (ii)~~\alpha_{j}(d\mu) = \alpha_{j}(d\nu)  & \mbox{   $j=0,1,\cdots,n-1$},\\
            (iii)~~\mu_{j}(d\mu) = \mu_{j}(d\nu)  & \mbox{  $j=0,1,\cdots,n$}.\\
            \end{array}  
    \end{eqnarray} 
\end{theorem}

\begin{proof}
    The recurrence (\ref{monic_recursion}) can be written in the matrix form
    $$
    \begin{pmatrix}
    \varphi_{j+1}^{R}(z) \\
    \varphi_{j+1}^{L,*}(z) 
    \end{pmatrix}= A^{R}(\alpha_{j},z)\begin{pmatrix}
    \varphi_{j}^{R}(z) \\
    \varphi_{j}^{L,*}(z) 
    \end{pmatrix},
    $$ 
    where
    $$
    A^{R}(\alpha,z) = \begin{pmatrix}
    z \rho^{-R} & -z \alpha \rho^{-L} \\
    - \alpha^{\dag} \rho^{-R}  & \rho^{-L}
    \end{pmatrix},
    $$
    and its inverse for $z \neq 0$,
    $$
    A^{-R}(\alpha,z) = \begin{pmatrix}
    z^{-1}  \rho^{-R} & \rho^{-R}\alpha\\
    z^{-1} \rho^{-L} \alpha^{\dag} & \rho^{-L}
    \end{pmatrix},
    $$
    which gives the inverse Szeg\H{o} relations 
    \begin{eqnarray} 
    \varphi_{j}^{R}(z) &=& z^{-1} \varphi_{j+1}^{R}(z) \rho_{j}^{-R} + z^{-1} \varphi_{j+1}^{L,*}(z) \rho_{j}^{-L} \alpha^{\dag}_j \label{monic_recursion_inv_1},\\
    \varphi_{j}^{L,*}(z) &=& \varphi_{j+1}^{R}(z) \rho_{j}^{-R} \alpha^{\dag}_j + \varphi_{j+1}^{L,*}(z) \rho_{j}^{-L}. \label{monic_recursion_inv_2}
    \end{eqnarray} 
    Then (\ref{eq:Phi_N}) implies $$\alpha_{n}(d\mu) = - \varphi_{n}^{\dag}(z;d\mu)\kappa_{n}^{L/2}= - \varphi_{n}^{\dag}(z;d\nu)\kappa_{n}^{L/2} = \alpha_{n}(d\nu),$$
    and thus by the inverse recursions (\ref{monic_recursion_inv_1}) and (\ref{monic_recursion_inv_2}), we have by iteration that (i) and (ii) hold.\\
    (i) implies (iii) because $\varphi_{j}(z;d\mu)$ and $\mu_{1}, \cdots, \mu_{j-1}$ determine $\mu_{j}(d\mu)$ via $$\int \Phi^{R}_{j}(z)d\mu=0.$$
\end{proof}
\noindent Turning to the weak convergence of $d\mu^{(n)}$:
\begin{theorem}
$d\mu^{(n)}$ is a probability measure on $\mathbb{T}$ for which (\ref{eq:BD_approx_bis}) holds. As $n \rightarrow \infty$, $d\mu^{(n)}\rightarrow d\mu$ weakly.
\end{theorem}
\begin{proof}
By using (iii) of Theorem (\ref{geronimus_thm}), for $j=0,1,\cdots,N$, we have $$\mu_{j}(d\mu^{(n)})=\mu_{j}(d\mu).$$
This equation and its conjugate imply that for any Laurent polynomial $f$ (i.e., a polynomial in $z$ and $z^{-1}$), we have
\begin{equation} \label{eq:coef_equality}
  \lim_{N \rightarrow \infty} \int f(e^{i\theta})d\mu^{(n)} = \int f(e^{i\theta})d\mu,
\end{equation}
because the left-hand side is equal to the right-hand side for large enough $N$. Since Laurent polynomials are dense in $C(\mathbb{T})$, equation (\ref{eq:coef_equality}) holds for all $f$. In other words, we have weak convergence.
\end{proof}


\section{The Szeg\H{o} Limit Theorem}

In this section, the main goal is to prove the Szeg\H{o} limit theorem for operator orthogonal polynomials. We emphasize that while operator-valued Szegő limit theorems have been obtained previously via Toeplitz operator factorization, the proof presented here is entirely based on operator-valued orthogonal polynomials and Schur complements.\\

Derevyagin, Holtz, Khrushchev, and Tyaglov (\cite{DerHK} Thm.28, \cite{Sim2} Thm 2.13.5) extended Szeg\H{o}'s theorem to the matrix version. Using their notation, det and tr representing the determinant and trace, respectively, the following limit holds:
\begin{equation} \label{KSz}
\lim_{n \rightarrow{ \infty}}\frac{\mathrm{det}(T_{n})}{\mathrm{det}(T_{n-1})}=\det \prod_0^{\infty} (I - {\alpha}_k {\alpha}_k^{\dag})
= \exp \int \hbox{tr} \ \log \mu'(\theta) d \theta /2 \pi .
\end{equation}
This is known as the Kolmogorov-Szeg\H{o} formula (see e.g. \cite{Bin1}, \S 4). Szeg\H{o}'s condition, which requires that the right-hand side of the equation be positive, holds if and only if
$$
\sum_0^{\infty} || {\alpha}_k^{\dag} {\alpha}_k || < \infty, 
$$
(extending early work of Delsarte, Genin and Kamp \cite{DelGK}).  \\
\noindent The {\it product theorem for determinants} in (\ref{KSz}) above is simple linear algebra in finitely many dimensions, and holds quite generally.\\

Before proving our results in the operator case, let us review some definitions related to infinite determinants \cite{Sim1}. For an operator $A$ that is trace class, the determinant $\mathrm{det}(I-A)$ is well-defined and can be written as
$$
\det(I-A)=\prod_{j=1}^{\infty}(1-\lambda_{j}),
$$ 
where $(\lambda_{j})_{j=1}^{\infty}$ are the eigenvalues of $A$.
On the other hand, if $A$ is a Hilbert-Schmidt operator, we define the second regularized determinant as
\[
\mathrm{det}_{2}(I-A)=\mathrm{det}[(I-A)e^{A}].
\] 
This is based on the observation that $(I-A)e^{A}-I$ is trace class. It is worth noting that when $A$ is trace class, we have
\[
\mathrm{det}_{2}(I-A)=\det(I-A)e^{tr(A)}.\]
Moreover, for two Hilbert-Schmidt operators $A$ and $B$, we have
\begin{equation} \label{eq:reg_det}
     \mathrm{det}_{2}(I-A)\mathrm{det}_{2}(I-B)=\mathrm{det}_{2}[(I-A)(I-B)]e^{tr(AB)}.
\end{equation} 
Let $\mathcal{W}_{1}(\mathcal{H})$, the Wiener algebra over the trace class operators on $\mathcal{H}$, be the set of all operator-valued functions $G$ on $\mathbb{T}$ of the form 
\begin{equation}\label{eq:W1_series_expansion}
G(z)=\sum_{n=-\infty}^{\infty}z^{n}G_{n}, ~~~ z \in \mathbb{T},
\end{equation} 
where $G_{n}$ is a trace class operator on $\mathcal{H}$ for each $n$ and 
\begin{equation}\label{eq:W1_space}
    \sum_{n=-\infty}^{\infty}||G_{n}||_{1} < \infty. 
\end{equation}

In the following theorem, we will prove the Szeg\H{o} Limit Theorem of the Szeg\H{o} theory of orthogonal polynomials and we assume $\mu_{0} = I$.

\begin{theorem}
Let $T^{R}_{n}$ be the Toeplitz operator matrix and $(\alpha_n)_{n}$ be the Verblunsky coefficients of $\mu$.  Then 
$$
\lim_{n \rightarrow{ \infty}}\frac{\mathrm{det}_{2}(T^{R}_{n})}{\mathrm{det}_{2}(T^{R}_{n-1})}=\det \prod_{k=0}^{\infty} (I - {\alpha}_k {\alpha}_k^{\dag}).
$$
\end{theorem}
\begin{proof}
    Using (\ref{eq:reg_det}), it follows that 
    \begin{eqnarray*}
        \mathrm{det}_{2}(T^{R}_{n}) &=& \mathrm{det}_{2}(T^{R}_{n-1})\mathrm{det}_{2}(I - \nu^{*}_{n} T^{-R}_{n}\nu_{n}) \exp({\mathrm{tr}(\nu^{*}_{n} T^{-R}_{n}\nu_{n})})\\
        &=&\mathrm{det}_{2}(T^{R}_{n-1})\mathrm{det}(I - \nu^{*}_{n} T^{-R}_{n}\nu_{n})\\
        &=& \mathrm{det}_{2}(T^{R}_{n-1}) \mathrm{det}(\kappa^{-R}_{n})=\mathrm{det}_{2}(T^{R}_{n-1}) \mathrm{det}(\rho_{0}^{R} \dots \rho_{n-1}^{R})^{2},     
    \end{eqnarray*}
    which leads to $$\mathrm{det}_{2}(T^{R}_{n})/\mathrm{det}_{2}(T^{R}_{n-1}) = \det \prod_{k=0}^{n-1} (I - {\alpha}_k {\alpha}_k^{\dag}). $$
    Since $T^{R}_{n}$ is definite positive, the coefficients $\alpha_{k}$ are Hilbert-Schmidt and strict contractions. This implies that ${\alpha}_k {\alpha}_k^{\dag}$ are trace class operators, and thus $\det (I - {\alpha}_k {\alpha}_k^{\dag})$ are well defined.
    Letting $n \rightarrow \infty$, the statement follows.
\end{proof}

As in Gohberg and Kaashoek \cite{GohK}, we assume in the theorem below that the measure $\mu'$ is in the Wiener algebra $\mathcal{W}_{1}(\mathcal{H})$.

\begin{theorem}
Assume $\mu'(z)=\sum_{j=-\infty}^{+\infty}\mu_{j}z^{j}\in I + \mathcal{W}_{1}(\mathcal{H})$ for $z \in \mathbb{T}$. Then 
 \begin{equation} \label{szego_thm}
      \lim_{n \rightarrow{ \infty}}\frac{\mathrm{det_{2}}(T^{R,L}_{n})}{\mathrm{det_{2}}(T^{R,L}_{n-1})}
 = \exp \frac{1}{2\pi}\int_{-\pi}^{\pi} \log(\det(\mu'(\theta))d\theta. 
 \end{equation}
\end{theorem}

\begin{proof}
 Write
    \[ 
    L =
    \begin{pmatrix}
    0 & \cdots & 0 & I\\
    0 &  \cdots & I & 0\\
    \vdots & I & \vdots & \vdots\\
    I & 0 & \cdots & 0
    \end{pmatrix},
    \qquad
    \nu_{n}=
    \begin{pmatrix}
    \mu_{-1}\\
    \mu_{-2}\\
    \vdots \\
    \mu_{-n}
    \end{pmatrix}.
    \]
    Then as $T^{R}_{n} = L T^{L}_{n} L$,
    \begin{eqnarray*}
    \Phi_{n}^{L,*} (z)
    &=& z^{n}(z^{-n}I-[I \ z^{-1}I \ \cdots \ z^{-n+1}I] \ T^{-L}_{n}{\nu}_{n}\\
    &=& I - [I~~zI~~ \cdots ~~ z^{n}I]L T^{-L}_{n} L \phi\\
    &=& I - [I~~zI~~ \cdots ~~ z^{n}I]T^{-R}_{n} \phi,\\
    &=& I + F_n(z),
    \end{eqnarray*}
    where $F_n(z) =-[I~~zI~~ \cdots ~~ z^{n}I]T^{-R}_{n} \phi$ and let $F(z) = \sum_{n=1}^{\infty}z^{n}F_{n}$.  \\
    
    Here, we can easily prove that the operator polynomial $F$ is trace class (since its coefficients are the sum of the product of two positive definite Hilbert-Schmidt operators). So $\det \Phi_{n}^{L,*} (z)$ and $\det (\Phi_{n}^{L,*})^{\dag} (z)$ are well-defined.\\
    
    Also, using the inequality (\ref{eq:W1_space}), we can conclude that the series on the right-hand side of (\ref{eq:W1_series_expansion}) converges in the trace class norm. Consequently, if we define $\mu'(\cdot) = I - \tilde{\mu}(\cdot)$ with $\tilde{\mu}$ belonging to the Wiener class $\mathcal{W}_{1}(\mathcal{H})$, then $\tilde{\mu}(z)$ is a trace class operator for all $z \in \mathbb{T}$. This means that $\det (\mu'(z))$ is well-defined for each $z \in \mathbb{T}$, and in particular, the expression $\det (\mu'(z))$ on the right-hand side of (\ref{szego_thm}) is well-defined.\\
    
    From (\ref{eq:BS_measure}) $\mu'^{
    (n)}=[\Phi_{n}^{L,*} \kappa_{n}^{-R}(\Phi_{n}^{L,*})^{\dag}]^{-1}$, and passing to the limit, we have $$\kappa_{\infty}^{R} = (I + F(z))^{\dag} \mu'(z) (I + F(z)).$$    
   It follows that  $$\log \det \kappa_{\infty}^{R}=
    M + 
    \frac{1}{2\pi}\int_{-\pi}^{\pi}\log \det (\mu'(\theta)) d\theta + \Bar{M},
    $$
    where 
    $$
    M = \frac{1}{2\pi}\int_{-\pi}^{\pi}\log \det (I + F(e^{it}))dt, \qquad
    \Bar{M} = \frac{1}{2\pi}\int_{-\pi}^{\pi}\log \det (I + F(e^{it}))^{\dag}dt.
    $$
    So $\det (I + F(\cdot)) $ is analytic on $|z|<1$, and $\det \Phi_{n}^{L,*}(\theta) $ is continuous and non-zero on $|z|\leq 1$. Combining, $\log \det (I + F(\cdot)) $ is analytic on $|z|<1$ and continuous on $|z|\leq 1$. Now Cauchy's theorem gives 
    $$
    M = \frac{1}{2\pi}\int_{-\pi}^{\pi}\log \det (I + F(e^{it}))dt = 0.
    $$
    Similarly for $\bar{M}=0$. Thus, we have the result.

\end{proof}

\section*{Declarations}
\ni \textbf{Funding} The authors received no specific funding for this research.\\

\ni \textbf{Conflict of interest} All authors declare that they have no conflict of interest

% \section{Complements}

% \ni 1. {\it Gaussian Regression Formula (GRF)} \\ 
% \i Much of what follows below is well illustrated by the {\it Gaussian Regression Formula (GRF)}.  If $X \sim N(\mu, \Si)$ is a Gaussian vector with  mean vector $\mu$ and covariance matrix $\Si = ({\sigma}_{ij})$, write the inverse covariance matrix ${\Si}^{-1}$ as $K = (k_{ij})$, the {\it concentration matrix}.  If $X, \ \mu, \ \Si, \ K$ are partitioned conformably,
% $$
% X = \begin{pmatrix}X_1 \cr
%                    X_2 \cr \end{pmatrix}, \quad
% \mu = \begin{pmatrix}{\mu}_1 \cr
%                    {\mu}_2 \cr \end{pmatrix}, \quad                   
% \Si = \begin{pmatrix}{\Si}_{11} & {\Si}_{12} \cr
%                      {\Si}_{21} & {\Si}_{22} \cr \end{pmatrix}, \quad
% K = \begin{pmatrix}{K}_{11} & {K}_{12} \cr
%                      {K}_{21} & {K}_{22} \cr \end{pmatrix},
% $$
% the conditional distribution of $X_2$ given $X_1 = x_1$ is
% $$
% X_2 \ | (X_1 = x_1) \sim N({\mu}_2 + {\Si}_{21} {\Si}_{11}^{-1} (x_1 - {\mu}_1, {\Si}_{22} - {\Si}_{21} {\Si}_{11}^{-1} {\Si}_{12}),                 \eqno(GRF)
% $$
% when ${\Si}_{11}^{-1}$ exists; see e.g. (\cite{BinF}, Th. 4.25).  For the general case, using the (Moore-Penrose) generalized inverse when ${\Si}$ is not invertible, see e.g. (\cite{Rao}, 8a.2.11).  Somewhat more simply,
% $$
% X_2 \ |(X_1 = x_1) \sim N({\mu}_2 - K_{22}^{-1} K_{21} (x_1 - {\mu}_1), K_{22}^{-1}).                                                              \eqno(GRF)                   
% $$
% Thus the regression of $X_2$ on $X_1$ is {\it linear}:
% $$
% \Ex [ X_2 \ | (X_1 = x_1)] = {\mu}_2 + {\Si}_{21} {\Si}_{11}^{-1} (x_1 - {\mu}_1)
% = {\mu}_2 - K_{22}^{-1} K_{21} (x_1 - {\mu}_1).
% $$
% \i The GRF has a long and remarkable history.  It originates in {\it Pearson's selection formula} of 1911-12 (selection being the term then used for what we now call regression); see e.g. (\cite{BinK}, \S 6.2). \\
% \i The GRF has been extended to the infinite-dimensional case.  See Hairer et al. (\cite{HaiS}, Lemma 4.4). \\

% \ni 2. {\it Schur complements} \\ 
% \i We have met the Schur complement already, under the notation $SC(.)$ (Issai Schur (1875-1941) in 1905; for background, see e.g. \cite{Zha}, \cite{BinK}).  If
% $$
% M =  \begin{pmatrix}P & Q \cr
%                      R & S \cr \end{pmatrix},
% $$
% the {\it Schur complement} of $P$ in $M$ is
% $$
% M/P := S - R P^{-1} Q.
% $$   
% The conditional variance matrix in $(GRF)$ above is thus the Schur complement of ${\Si}_{11}$ in $\Si$, $\Si/{\Si}_{11}$, also called the {\it partial covariance matrix} of $X_2$ given $X_1$.\\
% \i  For more on Schur complements in linear algebra, in the work of Schur, Aitken, Haynsworth and others, see (\cite{BinK}, \S 5.1).  For closely related inversion formulae (of Banachiewicz, Duncan and Bartlett, and the Sherman-Morrison-Woodbury formula), see (\cite{BinK}, \S 5.2).  These stem from the (surprisingly complicated) formula for the inverse of a partitioned matrix: in the notation above,
% $$
% M^{-1} =
% \begin{pmatrix}P^{-1} + P^{-1} Q (M/P)^{-1} R P^{-1} & -P^{-1} Q (M/P)^{-1} \cr
%                      -(M/P)^{-1} R P^{-1} & (M/P)^{-1}  \cr \end{pmatrix},
% $$
% when the relevant inverses exist; see e.g. (\cite{Dym}, (0.8) p.4), \cite{PunS}, \cite{TiaT}. \\
% \i Schur complements were extended to the operator case by Dritschel and Rovnyak \cite{DriR}.  See also Hairer et al. (\cite{HaiS}, Lemma 4.2). \\

% \ni 3. {\it Regression; independence and conditional independence} \\
% \i One sees (from the bivariate normal density and when it factorises) the classical result that for Gaussian vectors, two components $X_i, \ X_j$ are independent if and only if their correlation coefficient ${\sigma}_{ij} = 0$.  We note in passing the much more recent complement to this, {\it Dempster's theorem} of 1972 \cite{Dem}: two components $X_i, X_j$ of $X$ are {\it conditionally independent given all the rest} iff $k_{ij} = 0$.  For, taking $x_1$ for the random 2-vector $(X_i, X_j)$ and $x_2$ for the remainder that we condition on, the GRF shows that the conditional density $f_{1|2}$ of $x_1 | x_2$ has (conditional) covariance matrix $K_{11}^{-1}$.  So (repeating the argument just used) one has conditional independence iff $K_{11}^{-1}$ is diagonal, i.e (as $K_{11}$ is $2 \times 2$) $K_{11}$ diagonal, i.e. $k_{ij} = 0$ (\cite{BinK}, \S 6.3). \\
% \i For matrices $A, \ B$, linear forms $AX, \ BX$ are independent if and only if $A \Si B^T = 0$; see e.g. \cite{BinF}, Th. 4.16).  From this, one easily checks (\cite{BinF}, Th. 4.25, Proof) that $X_1$ and $X_2 - {\Si}_{21} {\Si}_{11}^{-1} X_1$ are independent, or:
% \begin{center}
% $X_1$ and $\Ex [ X_2 \ | \ X_1]$ are independent. 
% \end{center}
% This result extends to the infinite-dimensional case: for a version for Gaussian measures on locally convex topological vector spaces, see Bogachev (\cite{Bog}, 3.10.1). \\  

% \ni 4. {\it Shorted operators} \\
% \i The Schur complement occurs in a related context in the theory of {\it shorted operators}; see \cite{And}, \cite{AndT}, \cite{ButM}.  The term is suggested by the motivating physical context there, parallel connection of resistances in an electrical network. \\
% \ni 5. {\it Fej\'er-Riesz theorem} \\  
% \i Recall the classical {\it Fej\'er-Riesz theorem}: a trigonometric polynomial $Q(z) = \sum_{-n}^n Q_i z^i$ which is non-negative on the unit circle $\T$ can be factored as the square of an analytic trigonometric polynomial $P(z) = \sum_0^n P_i z^i$ (see e.g. (\cite{GreS}, 1.12 p.20), (\cite{Sim2}, 1.3.1 p.26)).  This was extended to the operator-valued case by Dritschel and Rovnyak \cite{DriR}: with $P^*$ the adjoint of $P$,
% $$
% Q(z) = P(z)^* P(z) \qquad (z \in \T).
% $$ 
% Their proof used operator-valued Schur complements (\cite{DriR}, \S 3).  \\

% \ni 6. {\it Riesz-Herglotz formula} \\ 
% \i Analytic functions $f: \D \to \C$ with $Re \ f(z) \geq 0 \ (z \in \D)$ normalised by $f(0) = 1$ are called {\it Carath\'eodory functions}.  The  Riesz-Herglotz formula is the representation of Carath\'eodory functions as
% $$
% f(z) =  
% \int_{\T} \ \Bigl( \frac{e^{i \theta} + z}{e^{i \theta} - z} \Bigr) \ d \mu (\theta),
% $$
% with $\mu$ a probability measure on $\T$ (\cite{Sim2}, 1.1 p.3).  Without the normalisation, one obtains the {\it Schur functions}, adds an $iC$ term on the right with $C$ real, and $\mu$ becomes a probability measure.  Operator versions of the Riesz-Herglotz formula are given by Dritschel and Rovnyak (\cite{DriR}; $\S$ 3, Method of Schur complements).  \\  

% \ni 7. {\it Operator form of Szeg\H{o}'s theorem} \\
% \i Rosenblum and Rovnyak (\cite{RosR}, \S 6.14) give an operator-valued version of Szeg\H{o}'s theorem.  With $F$ a non-negative operator-valued function satisfying a logarithmic integrability condition of Szeg\H{o} type, $F$ factorises as 
% $$
% F = G^* G,
% $$
% with $G$ an operator-valued (Hardy-space) outer function on $\T$.  See also Nikolskii (\cite{Nik1}, \S 4.8.8).  Recall the scalar form of Szeg\H{o}'s theorem, which involves the Szeg\H{o} function $h(z)$, an outer function, the `analytic square root' of the spectral density $w$ ($d\mu = w dm + d{\mu}_s$, with $\log w \in L^1$, to use the usual notation here, as in \cite{Sim2}, \cite{Bin1}).  \\
% \i The scalar form of Szeg\H{o}'s theorem is generalised by Dritschel and Rovnyak, to the matrix case (\cite{DriR}, Th. 4.7) and the operator case \cite{DriR}, Th. 4.5). \\
% \i For operator-valued spectral measures $\mu$, see Gamboa, Nagel and Rouault \cite{GamR}. \\
% \i There are many much more general related positivity results.  See e.g. Helton and Putinar \cite{HelP} for a survey of these, and applications to, e.g., optimization and control problems. \\

% \ni 8. {\it Operator form of Nehari's theorem} \\
% \i A {\it Hankel matrix} is one whose elements are constant on the backward diagonals, that is, one of the form $({\a}_{j+k})_{j,k = 0}^{\infty}$, for some sequence $\a = ({\a}_j)_0^{\infty}$ of complex numbers.  A {\it Hankel operator}   $\Gamma: {\ell}^2 \to {\ell}^2$ is obtained by extending the map  
% $$
% a = (a_j)_0^{\infty} \mapsto b = (b_j)_0^{\infty}, \qquad b_k := \sum_0^{\infty} {\a}_{j+k} a_j \quad (k \geq 0)
% $$
% from the (densely defined) sequences $a$ of finite support to ${\ell}^2$. \\
% \i  {\it Nehari's theorem} characterises the bounded (i.e. continuous, for linear operators) Hankel operators on ${\ell}^2$ as those with ${\a}_m$ here the non-negative Fourier coefficients $\hat \psi(m)$ of a bounded function $\psi \in L^{\infty} (\T)$.  Then $\Vert \Gamma \Vert$ is the infimum of $\Vert \psi \Vert_{\infty}$ over all such $\psi$.  See Nikolskii (\cite{Nik1}, \S 1.3, 1.4) for two proofs of Nehari's theorem, one by Riesz-Smirnov factorization, one by the stepwise extension method of Adamyan, Arov and Krein (`AAK', for brevity), below; see Peller \cite{Pel} for a monograph treatment of Hankel operators.\\
% \i The operator case is due to Page \cite{Pag} and AAK \cite{AdaAK2}, 1970/71; for a more recent treatment see Geronimo and Woerdeman \cite{GerW}. \\
% \i  For links between Nehari theory and prediction theory (below), see Yukio Kasahara and the first author \cite{KasB1}.  For the matrix case, see e.g. \cite{KasB2}, \cite{KasB3}. \\

% \ni 9. {\it The Nehari problem and rigidity} \\
% \i The term Nehari problem is used in two different but related ways. \\
% \i For the first, we turn to a reformulation of Nehari's theorem in terms of $H^2(\T)$ ($H^2$ for short), the Hardy class of functions of order 2 on the unit circle $\T$.  For $\phi \in L^2(\T)$ ($L^2$ below), define the {\it Hankel operator} $H_{\phi}$ from $H^2$ to $H^2_-$, its orthogonal complement in $L^2$, by 
% $$
% H_{\phi} f := P_ \phi f \qquad (f \in H^2),
% $$
% with $P_-$ the projection from $L^2$ to $H^2_-$.  Then $\phi$ is called a {\it symbol} of $H_{\phi}$ ({\it a} symbol, as symbols are far from unique).  Then Nehari's theorem may be reformulated as the equivalence, for $\phi \in L^2$, of: \\
% (i) the operator $H_{\phi}$ is bounded on $H^2$; \\
% (ii) $\phi$ has the same negative Fourier coefficients $\hat {\phi}_m$ ($m < 0$) as those, $\hat {\psi}_m$, of some function $\psi \in L^{\infty}$; \\
% (iii) $P_- \ \phi \in BMO$, the space of functions of bounded mean oscillation, \\
% and then
% $$
% \Vert H_{\phi} \Vert 
% = \inf \ \{ \Vert \psi \Vert_{L^{\infty}}: \hat \psi(m) = \hat \phi(m), \ m < 0 \}
% $$
% (Peller \cite{Pel}, Th. 1.3).  So, $H_{\phi}$ is bounded iff it has a bounded symbol, and as
% $$
% \Vert H_{\phi} \Vert 
% = \inf \{ \Vert \phi - f \Vert_{L^{\infty}}: f \in H^{\infty} \},
% $$
% the norm of $H_{\phi}$ is the distance from $\phi$ to $H^{\infty}$.  Hence the term {\it Nehari problem} for that of approximating to a bounded function (on $\T$) by bounded analytic functions.  See Peller \cite{Pel} (details below), Garnett (\cite{Gar}, IV.4,5) (without reference to Nehari, incidentally). \\
% \i For the second, one needs the concept of {\it rigidity} (the term is due to Sarason \cite{Sar}; cf. Poltaratski and Sarason \cite{PolS}).  A non-zero function $g \in H^1$ is {\it rigid} (or {\it strongly outer}) if it is {\it determined by its argument}, i.e. by $g/|g|$, to within a scale factor $c > 0$.  Then $h^2$ is an outer function in $H^1$, and $h^2/|h| = h/\bar h$ is called the {\it phase factor} of $h$.  The measure $\mu$ is determined by its Fourier coefficients, but with rigidity the {\it negative half} of these suffices. \\
% \i The second sense for `the Nehari problem' is: given a sequence $\g = ({\g}_n)_1^{\infty}$ of complex numbers, find $\phi$ in the unit ball of $L^{\infty}$ with
% $$
% {\g}_n = \int_{\T} e^{i n \theta} \phi dm \qq (n = 1,2, \cdots). \eqno (Neh)
% $$
% \i The case of non-uniqueness here -- the {\it indeterminate case} -- is the probabilistically interesting one.  In this case, when there is more than one (so infinitely many) solutions $\g$, $\g$ is called a {\it Nehari sequence} (below).  \\  

% \ni 10. {\it The Adamjan-Arov-Krein (AAK) parametrization} \\
% \i In the indeterminate case of the Nehari problem, the solutions were parametrised by Adamjan, Arov and Krein \cite{AdaAK1}.  They are of bilinear (fractional linear, M\"obius) type, parametrised by balls (the exceptional case of uniqueness occurs when the ball has radius 0).  For details see Peller \cite{Pel}: \S 5.1 (Th. 1.13 p.166, scalar case), and for the matrix- and operator-valued cases, \cite{AdaAK2}, and \cite{Pel}, \S 5.3 (Th. 3.5, p.177), \S 5.4 (Th. 4.16 p.212), \S 5.5 (Th. 5.8 p.225), \S 14.19 (Th. 19.1 p.628). \\

% \ni 11. {\it Nehari sequences and rigidity} \\
% \i In the indeterminate case, the AAK parametrisation describes the solution set as the outer functions $h \in H^2$ such that \\
% (i) $h$ has unit norm, \\
% (ii) $h^2$ is rigid, \\
% (iii) $h/\bar h$ solves the Nehari problem as in $(Neh)$, and one can take $h(0) > 0$. \\
% \i We now state the condition
% $$
% {\mu}_s = 0, \qq \log w \in L^1, \qq h^2 \ \hbox{rigid}.                \eqno(LM)
% $$
% Here $(LM)$ refers to work by Levinson and McKean \cite{LevM} in continuous time, for which this is the discrete-time analogue; we call $(LM)$ the case of `LM weights'.  Under $(LM)$, $\mu$ (and so $h$) is determined by its phase factor $h/\bar h$.  Restrict (as usual) to $\mu$ non-trivial (having infinitely many support points).  Then with $P_n$ the space of polynomials of degree $< n$ and ${\Phi}_n$ the monic orthogonal polynomials on $\T$ determined by $\mu$ (see e.g. \cite{Sim1}), $(LM)$ holds if and only if
% $$
% z^n H^2_-(\mu) \cap H^2(\mu) = P_n,
% $$
% equivalently,
% $$
% {\Phi}_n \ \perp z^n H^2_-(\mu) \cap H^2(\mu),
% $$
% for some ({\it equivalently, for all}) $n = 0,1,\cdots$ \ (\cite{KasB1}, Th. 3.1). \\
% \i For results in the matrix case, see \cite{KasB2}, \cite{KasB3}.  The operator case remains open. \\  

% \ni 12. {\it The moment problem} \\
% \i The indeterminate case of the Nehari problem suggests a comparison with the indeterminate case of the {\it moment problem}; both are strongly related to the theory of orthogonal polynomials.  For the extensive background here, see e.g. Khrushchev \cite{Bin4}. \\ 

% \ni 13. {\it Prediction theory} \\
% \i Our motivation here is probabilistic, prediction theory, in one, many or infinitely many dimensions \cite{Bin1}, \cite{Bin2}, \cite{Bin3}, \cite{BinM}.  The theory is at its most complete in the stationary case.  For the non-stationary case, see e.g. \cite{AlpDP}, \cite{DugBM}. \\ 
%%%%%%%%%%%%%%%%%%%%%%%%%%%%%%%%%%%%%%%%%%%%%%%%%%%%%%%%%%%%%
%%                  The Bibliography                       %%
%%                                                         %%
%%  imsart-???.bst  will be used to                        %%
%%  create a .BBL file for submission.                     %%
%%                                                         %%
%%  Note that the displayed Bibliography will not          %%
%%  necessarily be rendered by Latex exactly as specified  %%
%%  in the online Instructions for Authors.                %%
%%                                                         %%
%%  MR numbers will be added by VTeX.                      %%
%%                                                         %%
%%  Use \cite{...} to cite references in text.             %%
%%                                                         %%
%%%%%%%%%%%%%%%%%%%%%%%%%%%%%%%%%%%%%%%%%%%%%%%%%%%%%%%%%%%%%

%% if your bibliography is in bibtex format, uncomment commands:
%\bibliographystyle{imsart-number} % Style BST file (imsart-number.bst or imsart-nameyear.bst)
%\bibliography{bibliography}       % Bibliography file (usually '*.bib')

%% or include bibliography directly:
\begin{thebibliography}{9}
% \bibitem{AdaAK1} V. M. Adamyan, D. Z. Arov and M. G. Krein, Infinite Hankel matrices and generalized Carathéodory-Fejér and Schur problems. {\sl Funct. Anal. Appl.} {\bf2}(4), 1–17 (1968, in Russian). 
% \bibitem{AdaAK2} V. M. Adamyan, D. Z. Arov and M. G. Krein, Infinite block Hankel matrices and related extension problems. {\sl Izv. Akad. Nauk Armyan. SSR, Ser. Mat.} {\bf VI}(2–3), 87–112 (1971, in Russian).
% \bibitem{AlpDP} D. Alpay, A. Dijksma and Y. Peretz, Nonstationary analogs of the Herglotz representation theorem: the discrete case.  {\sl J. Funct. Anal.} {\bf 166} (1999), 85-129. 
% \bibitem{And} W. N. Anderson, Shorted operators.  {\sl SIAM J. Appl. Math.} {\bf 20} (1971), 520-525. 
% \bibitem{AndT} W. N. Anderson and G. E. Trapp, Shorted operators II.  {\sl SIAM J. Appl. Math.} {\bf 28} (1975), 60-71. 
\bibitem{Bin1} N. H. Bingham, Szeg\H{o}'s theorem and its probabilistic descendants.  {\sl Probability Surveys} {\bf 9} (2012), 287-324.
% \bibitem{Bin2} N. H. Bingham, Multivariate prediction and matrix Szeg\H{o} theory.  {\sl Probability Surveys} {\bf 9} (2012), 325-339. 
% \bibitem{Bin3} N. H. Bingham, Prediction theory for stationary functional time series.  {\sl Probability Surveys} {\bf 19} (2022), 160-184.
% \bibitem{Bin4} N. H. Bingham, The life, work and legacy of P. L. Chebyshev.  {\it P. L. Chebyshev - 200} (ed. A. N. Shiryaev).  {\sl Th. Prob. Appl.} {\bf 66} (4) (2022), 506-521. 
% \bibitem{BinF} N. H. Bingham and J. M. Fry, {\sl Regression: Linear models in statistics}.  Springer Undergraduate Mathematics Series (SUMS), Springer, 2010. 
% \bibitem{BinK} N. H. Bingham and W. J. Krzanowski, Linear algebra and multivariate analysis in statistics: development and interconnections in the twentieth century.  {\sl British Journal for the History of Mathematics} {\bf 37} no.1 (2022), 43-63.
% \bibitem{BinM} N. H. Bingham and Badr Missaoui, Aspects of prediction. {\sl J. Appl. Prob.} {\bf 51A} (2014), 189-201.
% \bibitem{Bog} V. I. Bogachev, {\sl Gaussian Measures}.  Math. Surv. Monog. {\bf 62}, Amer. Math. Soc., 1998.  
\bibitem{Bos} D. Bosq, Linear Processes in Function Spaces, {\sl Theory and Applications. Springer}, 2000.
\bibitem{BotS} A. B\"ottcher and B. Silbermann, Operator-Valued Szeg\H{o}-Widom Limit Theorems, Toeplitz Operators and Related Topics, {\sl Operator Theory: Advances and Applications}, vol. {\bf 71}, 1994, 33-53.
% \bibitem{ButM} C. A. Butler and T. D. Morley, A note on the shorted operator.  {\sl SIAM J. Matrix Anal. Appl.} {\bf 9} (1988), 147-155. 
\bibitem{Con} T. Constantinescu, Schur Parameters, Factorization and Dilation Problems. {\sl Birkh\"auser}, 1996. x+253 pp. 
\bibitem{DamPS} D. Damanik, A. Pushnitski and B. Simon, The analytic theory of matrix orthogonal polynomials. {\sl Surveys in Approximation Theory} {\bf 4} (2008), 1-85.
\bibitem{DerHK} M. Derevyagin, O. Holtz, S. Khrushchev and M. Tyaglov, Szeg\H{o}'s theorem for matrix orthogonal polynomials. {\sl J. Approx. Th.} {\bf 164} (2012), 1238-1261.
\bibitem{DelGK} P. Delsarte, Y. Genin and Y. G. Kamp, Orthogonal polynomial matrices on the unit circle.  {\sl IEEE Trans. Ciruits and Systems} CAS-{\bf 25} (1978), 149-160. 
% \bibitem{DriR} M. A. Dritschel and J. Rovnyak, The operator Fej\'er-Riesz theorem.  {\sl A glimpse at Hilbert space operators}, Paul R. Halmos in memoriam (ed. S. Axler, P. Rosenthal and D. Sarason), {\sl Opera. Th. Adv. Appl.} {\bf 207} (2010), 223-254.   
\bibitem{DugBM} M. Dugast, G. Bouleux and E. Marcon, Representation and characterization of non-stationary processes by dilation operators and induced shape space manifolds.  {\sl Entropy} {\bf 20} (2018), 23p. 
\bibitem{DunSa} N. Dunford and J. T. Schwartz, {\sl Linear operators, I: General theory}, Wiley, 1958. 
% \bibitem{DunSb} N. Dunford and J. T. Schwartz, {\sl Linear operators, II: Spectral theory, Self-adjoint operators in Hilbert space}, Wiley, 1963. 
% \bibitem{Dym} H. Dym, Krein's contributions to prediction theory. {\sl Operator Theory: Advances and Applications}, {\bf 118} (2000), 1-15, Birkh\"auser.
% \bibitem{GamR} F. Gamboa and A. Rouault. Operator-valued spectral measures and large deviations. {\sl J. Statist. Plann. Inference} {\bf 154} (2014), 72–86.
% \bibitem{Gar} J. B. Garnett, Bounded analytic functions. {\sl Pure and Applied Mathematics}, {\bf 96}. Academic Press, New York-London, 1981.
% \bibitem{GerW} J. S. Geronimo and H. J. Woerdeman, The operator-valued autoregressive filter problem and the suboptimal Nehari problem in two variables.  {\sl Int. Eq. Oper. Th.} {\bf 53} (2005), 343-361.
\bibitem{GohK} I. Gohberg and M.A. Kaashoek, Asymptotic Formulas of Szeg\H{o}-Kac-Achiezer Type, {\sl Asymptotic Analysis} {\bf 5} (1992), 187-22.
\bibitem{GorW} J. Gorniak and A. Weron, Aronszajn-Kolmogorov type theorems for positive definite kernels in locally convex spaces.  {\sl Studia Math.} {\bf 69} (1980), 45-56.
% \bibitem{GreS} U. Grenander and G. Szeg\H{o}, {\sl Toeplitz forms and their applications}.  {\sl University of California Press}, 1958.
% \bibitem{HaiS} M. Hairer, A. M. Stuart, J. Voss and P. Wiberg, Analysis of stochastic partial differential equations arising in path sampling.  I: The Gaussian case.  {\sl Comm. Math. Sci.} {\bf 3} (2005), 587-603. 
% \bibitem{HelP} J. W. Helton and M. Putinar, Positive polynomials in scalar and matrix variables, the spectral theorem and optimization.  {\sl Operator theory, structures matrices and dilations}, 229-306.  {\sl Theta Ser. Adv. Math.} {\bf 7}, 2007. 
% \bibitem{KasB1} Y. Kasahara and N. H. Bingham, Verblunsky coefficients and Nehari sequences. {\sl Trans. Amer. Math. Soc.} {\bf 366} (2014), 1363-1378.
% \bibitem{KasB2} Y. Kasahara and N. H. Bingham,  Coefficient stripping in the matricial Nehari problem.  {\sl J. Approximation Theory} {\bf 220} (2017), 1-11. 
% \bibitem{KasB3} Y. Kasahara and N. H. Bingham, Matricial Baxter's theorem with a Nehari sequence.  {\sl Mathematische Nachrichten} {\bf 291} (2018), 2590-2598. 
\bibitem{Kol} A. N. Kolmogorov, Stationary sequences in Hilbert space.  {\sl Bull. Moskov. Gos. Univ. Mat.} {\bf 2} (1941), 1-40 (in Russian; reprinted, {\sl Selected works of A. N. Kolmogorov, Vol. 2: Theory of probability and mathematical statistics}, Nauka, Moskva, 1986, 215-255).  
% \bibitem{LevM} N. Levinson and H. P. McKean, Weighted trigonometrical approximation on R1 with application to the germ field of a stationary Gaussian noise. {\sl Acta Math}. {\bf 112} (1964), 99-143.
% \bibitem{MarY} F. Marcell\'an and H. O. Yakhlef, Recent trends on analytic properties of matrix orthonormal polynomials, {\sl Electronic Transactions on Numerical Analysis}. {\bf 14}, (2002), 127-141.
\bibitem{ManS} V. Mandrekar and H. Salehi, The Square-integrability of Operator-valued Functions with Respect to non-negative Operator-valued Measure and The Kolmogorov Isomorphism Theorem, {\sl Indiana University Mathematics Journal}, vol. 20, {\bf 6}, 1970.
% \bibitem{Mas1} P. Masani, Orthogonally scattered measures.  {\sl Adv. Math.} {\bf 2} (1968), 61-117.
\bibitem{Mas2} P. Masani, Dilations as propagators of Hilbertian varieties, {\sl SIAM J. Math. Anal.} 9 (1978), 414-456.
% \bibitem{MiaS} A. G. Miamee and H. Salehi, The factorization problem for positive operator-valued functions in a Banach space and regularity of Banach space valued stationary stochastic processes.  {\sl Sanhy\=a A} {\bf 39} (1977), 211-222.
\bibitem{MiaS1} A. G. Miamee and H. Salehi, On the square root of positive $\mathcal{B}(\mathcal{X},\mathcal{X}^{*})$-valued function.  {\sl Journal of Multivariate Analysis} {\bf 7} (1977), 535-550.
% \bibitem{Mir} L. Miranian, Matrix valued orthogonal polynomials in the unit circle: some extensions of the clasical theory, {\sl Can. Math. Bulletin}, {\bf 52} (2009) 95-104.
% \bibitem{Nik1} N. K. Nikolskii, {\sl Treatise on the shift operator: Spectral function theory}. Grundl. Math. Wiss. {\bf 273}, Springer, 1986. 
\bibitem{Nik2} N. K. Nikolskii, {\sl Operators, functions and systems: an easy reading.  Volume 1: Hardy, Hankel and Toeplitz; Volume 2: Model operators and systems}.  Math. Surveys and Monographs {\bf 92, 93}, Amer. Math. Soc., 2002. 
% \bibitem{Pag} L. B. Page, Applications of the Sz.-Nagy and Foia\c{s} lifting theorem.  {\sl Indiana U. Math. J.} {\bf 20} (1970/71), 135-145. 
% \bibitem{Pel} V. V. Peller, {\sl Hankel operators and their applications}.  Springer, 2003.
% \bibitem{PolS} A. Poltoratski and D. Sarason, Aleksandrov-Clark measures, Recent advances in operator related function theory, 1–14, {\sl Contemp. Math.} {\bf 393}, {\sl Amer. Math. Soc.}, Providence, RI, 2006.
% \bibitem{PunS} S. Puntanen and G. P. H. Styan, Schur complements in probability and statistics.  In: Zhang, F. (ed.) The Schur Complement and Its Applications, pp. 163–226. Springer, US, Boston (2005). 
% \bibitem{Rao} C. R. Rao, {\sl Linear statistical inference and its applications}, 2nd ed., Wiley (1st ed. 1965). 
% \bibitem{RieN} F. Riesz and B. Sz.-Nagy, {\sl Le\c cons d'analyse fonctionelle}, 2nd ed., Akad\'emiai Kiad\'o, 1953 (1st ed. 1952). 
% \bibitem{RosR} M. Rosenblum and J. Rovnyak, {\sl Hardy classes and operator theory}.  Oxford University Press, 1985 (repr. Dover, 1997).
% \bibitem{Rud} W. Rudin, {\sl Functional analysis}, 2nd ed., McGraw-Hill, 1991 (1st ed. 1973).
% \bibitem{Sar} D. Sarason, Function theory on the unit circle, Notes for lectures given at a Conference at Virginia Polytechnic Institute and State University, 1978.
\bibitem{Sim1} B. Simon, Notes on infinite determinants of Hilbert space operators. {\sl Adv. Math.} {\bf{24}} (1977), 244-273.
\bibitem{Sim2} B. Simon, {\sl Orthogonal polynomials on the unit circle.  Part 1: Classical theory}. AMS Colloquium Publications 54.1, American Math. Soc., Providence RI, 2005.
% \bibitem{Sim3} Simon, B., The sharp form of the strong Szeg\H{o} theorem. {\sl Contemporary Math.} {\bf 387} (2005), 253-275, AMS, Providence RI.
\bibitem{Sze} G.~Szeg\H{o}, \emph{\"Uber den asymptotischen Ausdruck von Polynomen, die
durch eine Orthogonalit\"atseigenschaft definiert sind},
\emph{Math. Ann.} \textbf{86} (1922), 114--139.
\bibitem{SzNF} B. Sz.-Nagy, C. Foia\c s, H. Bercovici and L. K\'erchy, {\sl Harmonic analysis of operators on Hilbert space}, 2nd ed., Springer Universitext, 2010 (1st ed., Sz.-Nagy and Foia\c s, North-Holland, 1970).
\bibitem{Sto} M. H. Stone, {\sl Linear transformations in Hilbert space and their applications to analysis}.  Amer. Math. Soc. Colloq. Publ. {\bf XV}, AMS, 1932. 
% \bibitem{TiaT} Y. Tian and Y. Takane, The inverse of any two-by-two non-singular partitioned matrix and three matrix inverse completion problems.  {\sl Comput. Math. Appl.} {\bf 57} (2009), 1294-1304. 
\bibitem{Wer} A. Weron, Prediction theory in Banach spaces.  Proc. Winter Sch. Prob. Karpacz.  {\sl Lecture Notes in Math.} {\bf 472} (1975), 207-228.
\bibitem{WieM} N. Wiener and P. Masani, The prediction theory of multivariate stochastic processes II. The linear predictor. {\sl Acta Math.} {\bf 99} (1958), 93-137 (reprinted, Norbert Wiener, {\sl Collected Works III} (ed. P. Masani), MIT Press, 1981, p.204-248). 
% \bibitem{YMar} H. O. Yakhlef and F. Marcell\'an, Orthogonal Matrix Polynomials, Connection Between Recurrences on the Unit Circle and on a Finite Interval,  {\sl Approximation, Optimization and Mathematical Economics. Physica}, (2001).
\bibitem{Zha} F. Zhang, {\sl The Schur complement and its applications}.  {\sl Numerical Methods and Algorithms} {\bf 4}, Springer, 2005.
\end{thebibliography}

\end{document}
